\documentclass[11pt,reqno]{amsart}

\usepackage[T1]{fontenc}
\usepackage{lmodern}
\usepackage{microtype}
\usepackage{mathtools}
\usepackage{amssymb}
\usepackage{booktabs}
\usepackage{enumitem}
\usepackage{xcolor}
\usepackage[margin=1.08in]{geometry}
\usepackage[colorlinks=true,linkcolor=blue!55!black,citecolor=blue!55!black,
  urlcolor=blue!60!black]{hyperref}
\hypersetup{
  pdftitle={Universal Weighted-Laplacian Cones and the Maximal-Cusp Spectral Link},
  pdfauthor={Keio},
  pdfsubject={Weighted Laplacian inverse problems and small spectra of hyperbolic surfaces},
  pdfkeywords={weighted graph Laplacian, inverse eigenvalue problem, maximal-cusp spectral link, toric degeneration, strong spectral property, hyperbolic surface, semialgebraic complexity, sums of squares}
}

\newtheorem{theorem}{Theorem}[section]
\newtheorem{proposition}[theorem]{Proposition}
\newtheorem{corollary}[theorem]{Corollary}
\newtheorem{lemma}[theorem]{Lemma}
\theoremstyle{definition}
\newtheorem{definition}[theorem]{Definition}
\newtheorem{problem}[theorem]{Problem}
\theoremstyle{remark}
\newtheorem{remark}[theorem]{Remark}

\newcommand{\one}{\mathbf 1}
\newcommand{\R}{\mathbb R}
\newcommand{\Spec}{\operatorname{Spec}}
\newcommand{\cL}{\mathcal L}
\newcommand{\cR}{\mathcal R}
\newcommand{\SSPWL}{\ensuremath{\mathrm{SSPWL}}}
\newcommand{\Epos}{\mathcal E_{>0}}
\newcommand{\diag}{\operatorname{diag}}
\newcommand{\tr}{\operatorname{tr}}
\newcommand{\rank}{\operatorname{rank}}
\newcommand{\Sym}{\operatorname{Sym}}
\newcommand{\Crit}{\operatorname{Crit}}

\title[Universal Cones and Maximal-Cusp Spectral Links]{Universal
  Weighted-Laplacian Cones and the Maximal-Cusp Spectral Link}
\author{Keio}
\date{September 2026}

\subjclass[2020]{05C50, 15A18, 58J50, 30F60, 14P10, 68Q25}
\keywords{weighted graph Laplacian, inverse eigenvalue problem, maximal-cusp
  spectral link, stable graph, toric corner, strong spectral property,
  Ramanujan multigraph, finite-prefix nonrigidity, hyperbolic surface,
  semialgebraic complexity, sums of squares}

\begin{document}

\begin{abstract}
For each order we construct a scale-invariant cone of simple lacunary spectra
realized, at regular points with the strong spectral property, by every
connected weighted graph.  A toric blow-up of the rooted-forest coefficient
map gives the boundary scale--shape determinant $n/(n-1)$.  We also obtain the
sharp weighted-path eigenratio and an $O((\log n)/n)$ upper bound for its
normalized spectral volume.

The geometric part identifies the projectivized tangent link of the first
$2g-3$ positive eigenvalues at the spectrally maximal cusp of genus-$g$
moduli space.  The link is a finite union of normalized weighted-Laplacian
spectral images of complexity-one stable types; zero conductances encode
arbitrary relative pinching scales.  We make the required surface-uniform
Burger comparison explicit.  Outside a compact semialgebraic null skeleton,
every feasible simple direction lifts exactly at every sufficiently small
scale; weak-Arnold-hypothesis collision directions lift with multiplicity
preserved.

The link has sharp squared distance
$({g-3})/{(3(g-1)(2g-3))}$ from the equal-spectrum direction.  Hence a sharp
tangent cone is absent for $g\ge4$, whereas in genus $3$ every positive triple
ray, including repetitions, has a fixed-curvature exact lift.  A finite-size
bipartite Ramanujan construction supplies, in every genus, a full-dimensional
exact wedge of condition number below $18+12\sqrt2$; every target in it has
continuum many nonisometric realizations with the same entire sub-$1/4$
spectrum.  We further give an existential-real membership formulation,
graph-by-graph nonmembership certificates, and an exact-arithmetic
Frucht-fold certificate.  All surface conclusions beyond the tangent link
require a regular or essential graph preimage.  The finite-scale inverse
problem, including the residual collision skeleton and behavior near $1/4$,
remains open.
\end{abstract}

\maketitle

\section{Introduction}

The inverse eigenvalue problem for a weighted graph Laplacian asks which
lists
\[
  0=\lambda_1<\lambda_2\le\cdots\le\lambda_n
\]
can occur when the positive edge weights of a fixed connected graph are
allowed to vary.  Even for paths, the row-sum-zero constraint makes this
problem markedly different from the unconstrained symmetric inverse
eigenvalue problem.  Recent work of Fallat--Gupta--Lin
\cite{FGL2025} and Catral--Fallat--Gupta--Lin \cite{CFGL2026} developed
Jacobian and strong-property tools for this setting, while
Erchenko--Jakobson--Tsypin \cite{EJT2026} connected the realizable regions to
small eigenvalues of degenerating hyperbolic surfaces.
Mukherjee \cite{Mukherjee2026} proved approximate finite spectral
prescription for constant-negative-curvature manifolds when topology and the
curvature scale may vary; in dimension two that argument also uses the
classical collar graph limit.  Our surface results instead keep both the genus
and curvature $-1$ fixed and give exact realization on specified near-cusp
regions.

The usual direction is graph by graph: one searches for inequalities or
constructs a weighting adapted to the particular support.  Our first result
instead isolates a spectral region common to \emph{every} connected support.
For $0<\delta<1$ set
\begin{equation}\label{eq:lacunary-cone-intro}
 \cL_n(\delta)=
 \left\{0<\lambda_2<\cdots<\lambda_n:
  \frac{\lambda_i}{\lambda_{i+1}}<\delta
  \quad(2\le i\le n-1)\right\}.
\end{equation}

\begin{theorem}[Universal lacunary cone]\label{thm:intro-universal}
For every $n\ge2$ there is $\delta_n\in(0,1)$ such that
\[
  \cL_n(\delta_n)\subseteq\Epos(G)
\]
for every connected graph $G$ on $n$ vertices.  Every list in this cone has
a strict realization which is a regular point of the spectral map and has
the \SSPWL.
\end{theorem}

As of 4 September 2026, we are not aware of an earlier
support-independent open cone of this fixed spectral form lying in the strict
realization set of every connected graph on $n$ vertices.  Existing
strong-property results supply graph-dependent local spectral neighborhoods;
the common cone is the additional assertion here.

The mechanism behind Theorem~\ref{thm:intro-universal} is the main point.
Choose a spanning tree and order its edges so that the first $k$ edges form
a connected tree on $k+1$ vertices.  If successive edge weights are separated by independent small
ratios, the rooted-forest formula has a unique leading monomial in every
coefficient.  Dividing out those monomials produces an analytic map at a
toric corner.  Its shape Jacobian is diagonal there, and the determinant of
the full scale--shape Jacobian telescopes to
\[
  2\prod_{k=2}^{n-1}\left(1-\frac1{k^2}\right)=\frac{n}{n-1}.
\]
This supplies an exact local chart, including anisotropic lacunarity.  A
second implicit-function argument turns the zero weights on the remaining
edges of the graph into small positive weights without changing the
characteristic polynomial.

Our second theme is quantitative thinness.  For the path $P_n$, a classical
optimal-scaling argument of Forsythe--Straus \cite{ForsytheStraus1955}
specialized to the type-$A_{n-1}$ Cartan matrix gives
\begin{equation}\label{eq:path-ratio-intro}
  \frac{\lambda_n}{\lambda_2}
  \ge \cot^2\!\left(\frac{\pi}{2n}\right),
\end{equation}
with equality precisely for equal edge weights.  The inequality itself is
classical in this matrix-scaling guise; our application is to the spectral
volume $P(G)$ introduced in \cite{EJT2026}.  An exact beta integral for the
simplex region cut out by \eqref{eq:path-ratio-intro} yields an explicit
upper bound $P(P_n)\le U_n$, with
\[
 U_n=\frac{\pi^2}{4n}(\log n+\gamma_{\rm E})
 +O\!\left(\frac{(\log n)^2}{n^2}\right).
\]
Thus a universal nonempty cone can coexist with vanishing normalized volume.

We now turn to the geometric problem that motivated \cite{EJT2026}:
prescribing the first $2g-3$ eigenvalues of a closed curvature-$-1$ surface
of genus $g$.  Put
\[
 N=2g-2,\qquad m=N-1=2g-3,
\]
and, for a closed surface $X$, enumerate
\[
 0=\lambda_0(X)<\lambda_1(X)\le\lambda_2(X)\le\cdots
\]
and write
\[
 \Lambda_g(X)=(\lambda_1(X),\ldots,\lambda_m(X)),
 \qquad p(X)=\frac{\Lambda_g(X)}{\sum_{j=1}^m\lambda_j(X)}.
\]
The number $m$ is the full range below $1/4$ allowed by Otal--Rosas
\cite{OtalRosas2009}.  We call
\[
 \mathcal T_g=
 \left\{\lim_{k\to\infty}p(X_k):
       \lambda_m(X_k)\longrightarrow0\right\}
\]
the \emph{spectrally maximal-cusp tangent link}.  This terminology refers to
the maximal collapsing cluster; it does not require every Deligne--Mumford
node to open at one scale.

Let $\mathfrak G_g$ be the finite collection of complexity-one stable
genus-$g$ pseudographs described in Section~\ref{sec:tangent-link}.  For
$\Gamma\in\mathfrak G_g$, combine parallel nonloop cuff conductances, ignore
loops in the graph energy, and allow the effective conductances $w$ to vanish.
If $0=\mu_1\le\cdots\le\mu_N$ is the spectrum of $L_\Gamma(w)$, put
\[
 \psi_\Gamma(w)=
 \frac{(\mu_2,\ldots,\mu_N)}{\tr L_\Gamma(w)}.
\]

\begin{theorem}[Exact maximal-cusp tangent link]\label{thm:intro-link}
For every $g\ge2$,
\[
 \mathcal T_g=
 \bigcup_{\Gamma\in\mathfrak G_g}
 \left\{\psi_\Gamma(w):w\ge0,\ \tr L_\Gamma(w)>0\right\}.
\]
The link is compact and semialgebraic.  Zero conductances are essential: they
record multiscale pinching and may create additional zero coordinates in the
limiting direction.
\end{theorem}

On the simple positive part of the link, let $\mathcal U_g$ consist of the
directions having at least one normalized-regular graph preimage, and set
\[
 \mathfrak B_g=\mathcal T_g\setminus\mathcal U_g.
\]
Thus a simple point lies in $\mathfrak B_g$ precisely when every graph
preimage is critical; repeated and coordinate-wall directions are also placed
in the skeleton without assigning a derivative to the ordered eigenvalue map.

\begin{theorem}[Generic exactness]\label{thm:intro-generic}
The set $\mathcal U_g$ is relatively open and semialgebraic in the strict
ordered simplex (rather than merely relatively open in $\mathcal T_g$), while
$\mathfrak B_g$ is compact and semialgebraic with
\[
 \dim\mathfrak B_g\le m-2,
 \qquad \mathcal H^{m-1}(\mathfrak B_g)=0.
\]
Here $\mathfrak B_2=\varnothing$ and $\dim\varnothing=-\infty$.
For every $p\in\mathcal U_g$ there is $s_0(p)>0$ such that every $sp$,
$0<s<s_0(p)$, is exactly the first $m$ positive-eigenvalue list of a closed
genus-$g$ hyperbolic surface.  If $p\notin\mathcal T_g$, then no $sp$ is
realizable for all sufficiently small $s>0$.
\end{theorem}

The regularity hypothesis is sufficient, not necessary.  A form-valued
version of the Colbois--Colin de Verdi\`ere essential-family argument
\cite{ColboisCdV1988} also lifts collision directions admitting the weak
Arnold hypothesis (WAH); the strong Arnold hypothesis (SAH), and in particular
the \SSPWL, implies WAH.  Directions not settled by either mechanism remain
confined to the lower-dimensional skeleton.

\begin{theorem}[Conditioned wedge and maximal-prefix nonrigidity]
\label{thm:intro-conditioned}
For every $g\ge2$ there are a nonempty open set
$V_g\subset\{0<\eta_1<\cdots<\eta_{2g-3}\}$ and $s_g>0$ such that, for
every $\eta\in V_g$ and $0<s<s_g$,
\[
 \frac{\eta_{2g-3}}{\eta_1}<18+12\sqrt2,
\]
and continuum many pairwise nonisometric closed genus-$g$ surfaces have
$s\eta$ as their entire sub-$1/4$ spectrum.  For each fixed $s$, almost
every target in $sV_g$ has a regular marked realization with local isolevel
dimension $4g-3$.
\end{theorem}

The next eigenvalue is always larger than $1/4$, so the outside gap can be
made arbitrarily large while retaining the displayed uniform internal
bound.  In addition, for large $g$ one common exact maximal prefix appears
in at least $\exp\{g\log g-O(g)\}$ canonically distinguishable all-pants
thin regions.  The word \emph{Ramanujan} refers here only to the finite
cubic bipartite dual multigraph supplied by \cite{MSS2018}; no automorphic
or Selberg-type assertion is involved.

The topology of a stable genus-$g$ graph forces a sharp quantitative gap.
Write
\[
 b_m=(1/m,\ldots,1/m).
\]

\begin{theorem}[Sharp cusp gap and the genus-$3/4$ transition]
\label{thm:intro-cusp-gap}
For every $g\ge3$,
\[
 \lim_{\varepsilon\downarrow0}
 \inf_{\lambda_{2g-3}(X)\le\varepsilon}
 \|p(X)-b_m\|_2^2
 =\frac{g-3}{3(g-1)(2g-3)}.
\]
Consequently, for every $g\ge4$ and every
$r<\sqrt{(g-3)/(3(g-1)(2g-3))}$, the radius-$r$ angular cone about the
equal-spectrum ray is absent sufficiently near the origin; the displayed
radius is sharp at the tangent level.  In contrast, for $g=3$ and every
$0<a\le b\le c$, there is $s_0(a,b,c)>0$ such that every $s(a,b,c)$,
$0<s<s_0(a,b,c)$, is exactly the first three positive eigenvalues of a
closed genus-$3$ hyperbolic surface, with multiplicities preserved.
\end{theorem}

If
\[
 \kappa_g^*=\inf_{p\in\mathcal U_g}\frac{p_m}{p_1},
 \qquad N=2g-2,
\]
then the sharp variance theorem and the conditioned wedge give, for $g\ge3$,
\[
 \frac{5N-8+4\sqrt{(N-1)(N-4)}}{3N}
 \le\kappa_g^*<18+12\sqrt2.
\]
The lower bound tends to $3$; this bracket does not determine
$\kappa_g^*$.

Colin de Verdi\`ere's property $A_4$ gives the corresponding unnormalized
constant-curvature finite-spectrum prescription in genus three, with four
eigenvalues including the zero mode and with repetitions allowed
\cite[Theorem~1.3 and \S2a]{CdV1987}.  That theorem does not fix the
curvature normalization.  The statement here keeps curvature equal to $-1$
and realizes every sufficiently small point on each ray by a direct $K_4$
essential-family argument; its role is to identify the zero-gap endpoint of
the maximal-cusp link theorem.

The universal cone remains a concrete regular chart inside this larger
picture.

\begin{theorem}[Exact lacunary prescription on a fixed-genus surface]
\label{thm:intro-surface}
For every $g\ge2$ there is $\varepsilon_g>0$ with the following property.
If
\[
 0<\rho_1<\cdots<\rho_{2g-3},\qquad
 \frac{\rho_j}{\rho_{j+1}}<\varepsilon_g
 \quad(1\le j\le2g-4),
\]
then there is $s_0=s_0(\rho)>0$ such that, for every $0<s<s_0$, a closed
hyperbolic surface $X_s$ of genus $g$ has
\[
  \lambda_j(X_s)=s\rho_j,
  \qquad 1\le j\le2g-3,
 \qquad \lambda_{2g-2}(X_s)>\frac14.
\]
\end{theorem}

The final layer is effective in the sense of real algebraic geometry.  For a
rational direction $p$, tangent-link membership has a polynomial-size
existential-real encoding.  For fixed genus, standard fixed-dimensional
quantifier-elimination bounds give polynomial bit complexity.  Off the link,
strict positivity of a trace-moment error polynomial on each compact weight
simplex gives a finite Putinar--Lasserre sums-of-squares certificate.  These
are upper bounds and certificates only: no hardness result and no uniform SOS
degree bound is asserted.

The logical scope is worth recording explicitly.  The tangent-link theorem is
a complete asymptotic classification of projectivized maximal collapsing
clusters.  Exact statements at positive scale require either a regular graph
preimage or an essential form-valued family.  We do not solve the full
finite-scale problem of \cite[Question~1]{EJT2026}, classify the residual
collision skeleton, or control spectral continuation through the threshold
$1/4$.

The paper is organized as follows.  Sections~\ref{sec:prelim}--
\ref{sec:universal} prove the universal graph theorem and record a sharp
generic algebraic-degree bound for trees.  Sections~\ref{sec:path}--
\ref{sec:volume} give the independent path obstruction and volume consequence.
Sections~\ref{sec:surface}--\ref{sec:full-index} establish the exact
graph-to-surface transfer.  Sections~\ref{sec:tangent-link}--
\ref{sec:genus-three} prove the tangent-link, generic-exactness, conditioned
maximal-cluster, finite-prefix nonrigidity, essential-collision, sharp-gap,
and genus-three theorems.  Section~\ref{sec:complexity}
develops the algebraic decision and SOS certificates, and
Section~\ref{sec:residual} isolates the unresolved singular directions and
the global continuation problem.

\section{Weighted Laplacians and spectral coordinates}\label{sec:prelim}

Let $G=(V,E)$ be a simple graph on $V=[n]$.  Give each edge $e=uv$ an
orientation and put $b_e=\mathbf e_u-\mathbf e_v$.  For
$w\in\R_{>0}^{E}$ its strictly weighted Laplacian is
\begin{equation}\label{eq:weighted-L}
 L_G(w)=\sum_{e\in E}w_eb_eb_e^{\mathsf T}.
\end{equation}
The orientation is immaterial.  We denote by $\Epos(G)$ the set of ordered
positive spectra realized by such weights.  When $G$ is connected,
$L_G(w)$ is positive semidefinite with kernel $\R\one$.

Put $m=n-1$ and write the characteristic polynomial in the form
\begin{equation}\label{eq:q-def}
 \det(tI_n+L_G(w))
 =t\prod_{i=2}^n(t+\lambda_i)
 =\sum_{k=0}^{m}q_k(w)t^{n-k},
 \qquad q_0=1.
\end{equation}
Thus $q_k=e_k(\lambda_2,\ldots,\lambda_n)$.  A realization is
\emph{coefficient-regular} if the map
\[
 w\longmapsto(q_1(w),\ldots,q_m(w))
\]
has rank $m$.  For a simple spectrum this is equivalent to regularity of the
ordered spectral map.

We will repeatedly use the following elementary coordinate system on the
positive coefficient orthant.  For $2\le k\le m$ define
\begin{equation}\label{eq:theta-def}
 \vartheta_k(q)=\frac{q_kq_{k-2}}{q_{k-1}^2}.
\end{equation}
Then
\begin{equation}\label{eq:q-recursion}
 q_k=\vartheta_k\frac{q_{k-1}^2}{q_{k-2}},
\end{equation}
so $q\mapsto(q_1,\vartheta_2,\ldots,\vartheta_m)$ is a diffeomorphism of
the positive orthant.

Let $u_2,\ldots,u_n$ be orthonormal eigenvectors for a simple nonzero
spectrum.  The eigenvalue derivative is
\begin{equation}\label{eq:eigen-derivative}
 \frac{\partial\lambda_i}{\partial w_e}
 =(b_e^{\mathsf T}u_i)^2.
\end{equation}
The matrix on the right, with rows indexed by $i=2,\ldots,n$, is the
spectral Jacobian with respect to the edge weights.  Since the Jacobian from
the simple eigenvalues to their elementary symmetric functions is
nonsingular, coefficient regularity is equivalent to full row rank of this
spectral Jacobian.  It is not the verification matrix of
\cite{CFGL2026}; rather, the implication from this regularity to the
\SSPWL follows from \cite[Theorem~3.8]{CFGL2026}.

\section{Spanning forests and toric corner charts}\label{sec:toric}

The coefficient blow-up rests on the rooted-forest formula.  For a forest
$F$, let $\pi(F)$ denote all components of the spanning forest $(V,F)$,
including isolated vertices.

\begin{lemma}[Rooted-forest coefficients]\label{lem:forest}
For $0\le k\le n-1$,
\begin{equation}\label{eq:forest-formula}
 q_k(w)=
 \sum_{\substack{F\subseteq E(G),\ |F|=k\\F\text{ a forest}}}
 \left(\prod_{C\in\pi(F)}|C|\right)\prod_{e\in F}w_e.
\end{equation}
\end{lemma}

\begin{proof}
Let $B$ be an oriented incidence matrix and $W=\diag(w_e)$, so
$L=BWB^{\mathsf T}$.  The coefficient $q_k$ is the sum of the principal
$k$-minors of $L$.  For $S\subseteq V$, $|S|=k$, Cauchy--Binet gives
\[
 \det L[S,S]=
 \sum_{\substack{F\subseteq E(G)\\|F|=k}}
 \det B[S,F]^2\prod_{e\in F}w_e.
\]
The determinant is nonzero precisely when $F$ is a forest and every
component contains exactly one vertex of $V\setminus S$; then its square is
one.  For fixed $F$, choosing $V\setminus S$ is therefore the same as
choosing one root in each component, in
$\prod_{C\in\pi(F)}|C|$ ways.  Summing over $S$ proves
\eqref{eq:forest-formula}.  This is the rooted-forest specialization of the
all-minors matrix--tree theorem; see \cite{Chaiken1982}.
\end{proof}

Let $T$ be a tree on $n$ vertices.  Fix an ordering
\begin{equation}\label{eq:prefix-order}
 E(T)=\{e_1,\ldots,e_m\}
\end{equation}
such that the first $k$ edges form a connected tree on $k+1$ vertices for
each $k$.  Indeed, root $T$, list the nonroot vertices in a
parent-before-child order, and let $e_k$ join the $k$th listed vertex to its
parent.  Equivalently, the reverse order is a pendant-edge deletion order.

Introduce the toric weight coordinates
\begin{equation}\label{eq:toric-weights}
 w_1=s,\qquad w_j=sx_1\cdots x_{j-1}\quad(2\le j\le m),
\end{equation}
where $s>0$.  For $I\subseteq[m]$ and $1\le r\le m-1$, put
\[
 \nu_r(I)=|\{i\in I:i>r\}|.
\]
If $|I|=k$, then
\begin{equation}\label{eq:valuation-ineq}
 \nu_r(I)\ge(k-r)_+=\nu_r([k]).
\end{equation}
Equality for every $r$ holds only for $I=[k]$.  Consequently, after
\eqref{eq:toric-weights}, every monomial in $q_k$ is divisible by
\[
 s^kX_k(x),\qquad X_k(x)=\prod_{r=1}^{k-1}x_r^{k-r}.
\]
The prefix forest $\{e_1,\ldots,e_k\}$ has one component of size $k+1$
and all other vertices isolated.  Lemma~\ref{lem:forest} and
\eqref{eq:valuation-ineq} therefore give the exact factorization
\begin{equation}\label{eq:tree-factorization}
 q_k^T(s,x)=(k+1)s^kX_k(x)A_{T,k}(x),
 \qquad A_{T,k}(0)=1,
\end{equation}
where $A_{T,k}$ is a polynomial depending on this fixed prefix ordering.
We set $X_0=A_{T,0}=1$.

Combining \eqref{eq:tree-factorization} with \eqref{eq:theta-def} yields
\begin{equation}\label{eq:tree-theta}
 \vartheta_k^T(x)=
 \gamma_kx_{k-1}
 \frac{A_{T,k}(x)A_{T,k-2}(x)}{A_{T,k-1}(x)^2},
 \qquad
 \gamma_k=\frac{(k-1)(k+1)}{k^2}=1-\frac1{k^2}.
\end{equation}
Thus the blown-up shape map
\[
 \Phi_T(x)=(\vartheta_2^T(x),\ldots,\vartheta_m^T(x))
\]
extends analytically across the corner and
\begin{equation}\label{eq:tree-jacobian}
 D\Phi_T(0)=\diag(\gamma_2,\ldots,\gamma_m).
\end{equation}
For the analytically extended blown-up map, including the scale coordinate
gives the exact boundary determinant
\begin{equation}\label{eq:telescoping-det}
 \det D_{(s,x)}(q_1,\vartheta_2,\ldots,\vartheta_m)\big|_{(s,x)=(s,0)}
 =2\prod_{k=2}^m\left(1-\frac1{k^2}\right)
 =\frac{m+1}{m}=\frac{n}{n-1}.
\end{equation}
Here the derivative in $s$ is $2$ at the corner; the determinant is
independent of the positive boundary scale.  Formula
\eqref{eq:telescoping-det} is the transversality mechanism of the paper.

\section{A universal lacunary realization cone}\label{sec:universal}

Write a target list in decreasing order as
\[
 \ell_1=\lambda_n,\quad \ell_2=\lambda_{n-1},\quad\ldots,\quad
 \ell_m=\lambda_2
\]
and set
\begin{equation}\label{eq:spectral-toric}
 \ell_1=\tau,\qquad
 \ell_j=\tau z_1\cdots z_{j-1}\quad(2\le j\le m).
\end{equation}
Then $z_j=\ell_{j+1}/\ell_j$, so $0<z_j<\delta$ is exactly the
lacunarity condition.  If $p_k=e_k(\ell_1,\ldots,\ell_m)$, the same monomial
argument as above, now with every subset coefficient equal to one, gives
\begin{equation}\label{eq:spectral-factorization}
 p_k(\tau,z)=\tau^kX_k(z)B_k(z),\qquad B_k(0)=1.
\end{equation}
We set $B_0=1$.
Hence
\begin{equation}\label{eq:spectral-theta}
 \vartheta_k^{\rm sp}(z)=
 z_{k-1}\frac{B_k(z)B_{k-2}(z)}{B_{k-1}(z)^2}.
\end{equation}
The map
\[
 \Phi_{\rm sp}(z)=
 (\vartheta_2^{\rm sp}(z),\ldots,\vartheta_m^{\rm sp}(z))
\]
is analytic near zero, with
$\Phi_{\rm sp}(0)=0$ and $D\Phi_{\rm sp}(0)=I_{m-1}$.

\begin{proposition}[Tree lacunary chart]\label{prop:tree-chart}
For every tree $T$ on $n$ vertices there is $\delta_T>0$ such that every
list in $\cL_n(\delta_T)$ is realized by strictly positive weights on $T$.
The realization is coefficient-regular.
\end{proposition}

\begin{proof}
By \eqref{eq:tree-jacobian} and the analytic inverse-function theorem,
after shrinking neighborhoods, $\Phi_T$ is a diffeomorphism
$U_T\to V_T$ about zero.  We may suppose that all factors
\[
 H_{T,k}(x)=
 \frac{A_{T,k}(x)A_{T,k-2}(x)}{A_{T,k-1}(x)^2}
\]
are positive on $U_T$.  Choose $\delta_T>0$ so small that the analytically
extended map satisfies
\[
 \Phi_{\rm sp}\big([0,\delta_T]^{m-1}\big)\subseteq V_T.
\]
For a target \eqref{eq:spectral-toric}, set
\[
 x=\Phi_T^{-1}(\Phi_{\rm sp}(z)).
\]
This inverse lies in the positive orthant.  Indeed,
\eqref{eq:tree-theta} gives coordinate by coordinate
\[
 \vartheta_k^{\rm sp}(z)=\gamma_kx_{k-1}H_{T,k}(x),
\]
whose left side and all factors except $x_{k-1}$ are positive.

It remains to match scale.  Since
\[
 q_1^T=2sA_{T,1}(x),\qquad p_1=\tau B_1(z),
\]
choose
\begin{equation}\label{eq:scale-choice}
 s=\frac{\tau B_1(z)}{2A_{T,1}(x)}>0.
\end{equation}
Then $q_1^T=p_1$ and all the $\vartheta_k$ agree.  The recursion
\eqref{eq:q-recursion}, starting from $q_0=p_0=1$, proves
$q_k^T=p_k$ for all $k$.  Thus the characteristic polynomials, and hence
the spectra, agree exactly.

Finally, $\Phi_T$ is a local diffeomorphism and
$\partial q_1^T/\partial s=2A_{T,1}(x)>0$.  Both
$(s,x)\mapsto w$ and
$q\mapsto(q_1,\vartheta_2,\ldots,\vartheta_m)$ are diffeomorphisms on the
positive orthant.  The coefficient Jacobian with respect to the $m$ tree
weights is therefore nonsingular.
\end{proof}

The chart also records the leading weights.  As $z\to0$,
\begin{equation}\label{eq:weight-asymptotic}
 x_{k-1}=\frac{k^2}{k^2-1}z_{k-1}(1+O(\|z\|)),
 \qquad
 w_j=\frac{j}{j+1}\ell_j(1+O(\|z\|)).
\end{equation}
The second identity uses
$\prod_{k=2}^j k^2/(k^2-1)=2j/(j+1)$ and
$s=\ell_1/2\,(1+O(\|z\|))$.

To pass from a spanning tree to the full graph we need an exact, not merely
limiting, perturbation.

\begin{lemma}[Regular spanning-subgraph lift]\label{lem:supergraph}
Let $H$ be a connected spanning subgraph of $G$.  Suppose a simple spectrum
has an $H$-realization at which the coefficient map has full rank with
respect to $n-1$ selected edge weights of $H$.  Then the same spectrum has a
strictly positive, coefficient-regular realization on $G$.
\end{lemma}

\begin{proof}
Hold every unselected edge weight of $H$ fixed at its original positive
value.  Separate the selected weights as $u\in\R^{n-1}$ and the new edge
weights as $v$.  Set
\[
 F(u,v)=(q_1(L_G(u,v)),\ldots,q_m(L_G(u,v))).
\]
At the original realization $(u^0,0)$, the derivative $D_uF(u^0,0)$ is
invertible.  The implicit-function theorem supplies $u=u(v)$ near zero with
$F(u(v),v)=F(u^0,0)$.  Choose every coordinate of $v$ positive and
sufficiently small.  Continuity gives $u(v)>0$ and preserves the
invertibility of $D_uF$.  The new Laplacian has support exactly $G$ and the
same characteristic polynomial.
\end{proof}

\begin{proof}[Proof of Theorem~\ref{thm:intro-universal}]
For $n=2$, take $\delta_2=1/2$; the only connected graph is $K_2$, and
$w=\lambda_2/2$ is a regular realization, hence has the \SSPWL.  For
$n\ge3$, choose one prefix order for each
labeled tree $T$ on $n$ vertices, and let $\delta_T$ be supplied by
Proposition~\ref{prop:tree-chart}.  There are finitely many such trees, so
\[
 \delta_n=\min\left\{\frac12,\min_T\delta_T\right\}>0.
\]
Given a connected $G$, choose a spanning tree $T\subseteq G$.
Proposition~\ref{prop:tree-chart} gives a regular strict realization on $T$,
and Lemma~\ref{lem:supergraph} lifts it to a regular strict realization on
$G$.  The selected tree-weight columns form a nonsingular $m\times m$ minor
of the full coefficient Jacobian on $G$.  Since the prescribed spectrum is
simple, \cite[Theorem~3.8]{CFGL2026} gives the \SSPWL.
\end{proof}

\begin{remark}[What is and is not uniform]
The number $\delta_n$ depends only on $n$, not on $G$, the overall scale, or
the anisotropy among the successive ratios.  The proof is qualitative: it
does not optimize $\delta_n$.  It also requires strict inequalities and a
simple target spectrum.  Boundary spectra with repeated or zero eigenvalues
are outside this chart.
\end{remark}

The same regular point also controls the generic algebraic complexity of the
tree inverse problem.  The following degree is taken over $\mathbb C$, with
edge labels retained and isolated solutions counted with multiplicity.

\begin{proposition}[Sharp universal generic-degree bound for trees]
\label{prop:tree-degree}
Let $T$ be a labeled tree on $n$ vertices and put $m=n-1$.  On the affine
trace hyperplane
\[
 H_T=\{w\in\mathbb C^m:q_1(w)=1\},
\]
the coefficient map
\[
 \mathcal C_T:H_T\longrightarrow\mathbb C^{m-1},\qquad
 \mathcal C_T(w)=(q_2(w),\ldots,q_m(w)),
\]
is dominant and generically finite, and
\[
 \deg\mathcal C_T\le 2\cdot3\cdots m=m!.
\]
This bound is sharp: for the labeled star $K_{1,m}$ one has
$\deg\mathcal C_{K_{1,m}}=m!$.
\end{proposition}

\begin{proof}
Proposition~\ref{prop:tree-chart} supplies a positive point at which the full
coefficient Jacobian is nonsingular.  Homogeneity separates the radial trace
direction, so the derivative of $\mathcal C_T$ restricted to $T_wH_T$ is
nonsingular after rescaling that point to $q_1=1$.  Hence $\mathcal C_T$ is
dominant and generically finite.  Eliminate one weight with the linear
equation $q_1=1$.  A generic fiber is then cut out by equations of degrees
$2,3,\ldots,m$, so affine B\'ezout gives the stated upper bound
\cite{CoxLittleOShea2005}.

For the star, every $k$-edge forest has one nontrivial component of order
$k+1$.  The rooted-forest formula therefore gives
\[
 q_k(w)=(k+1)e_k(w_1,\ldots,w_m).
\]
Together with the fixed value of $e_1$, a generic target determines the monic
polynomial having $w_1,\ldots,w_m$ as its roots.  Distinct roots give exactly
the $m!$ labeled permutations.  Thus equality holds.  This is a complex
generic degree statement; a particular real fiber can be smaller or
positive-dimensional on the exceptional algebraic set.
\end{proof}

\section{A sharp all-order obstruction for weighted paths}\label{sec:path}

Let $P_n$ have edges $e_i=\{i,i+1\}$, $1\le i\le n-1$, and let
$D\in\R^{n\times(n-1)}$ be the oriented incidence matrix whose $i$th column
is $\mathbf e_i-\mathbf e_{i+1}$.  For positive weights
$w_1,\ldots,w_{n-1}$ put $W=\diag(w_1,\ldots,w_{n-1})$.  The nonzero
eigenvalues of $L_w=DWD^{\mathsf T}$ are the eigenvalues of
\begin{equation}\label{eq:cartan-congruence}
 A_w=W^{1/2}C_{n-1}W^{1/2},\qquad
 C_{n-1}=D^{\mathsf T}D=
 \begin{pmatrix}
 2&-1\\
 -1&2&\ddots\\
 &\ddots&\ddots&-1\\
 &&-1&2
 \end{pmatrix}.
\end{equation}

\begin{proposition}[Sharp path eigenratio]\label{prop:path-ratio}
For every $n\ge2$ and every positive weighting of $P_n$,
\begin{equation}\label{eq:path-ratio}
 \frac{\lambda_n(L_w)}{\lambda_2(L_w)}
 \ge\frac{1+\cos(\pi/n)}{1-\cos(\pi/n)}
 =\cot^2\!\left(\frac{\pi}{2n}\right).
\end{equation}
Equality holds if and only if $w_1=\cdots=w_{n-1}$; for $n=2$ this
condition is vacuous.
\end{proposition}

\begin{proof}
Set
\[
 s_i=\sin\frac{i\pi}{n},\qquad
 t_i=(-1)^{i+1}s_i,\qquad 1\le i\le n-1.
\]
The two endpoint eigenpairs of $C_{n-1}$ are
\[
 C_{n-1}s=\alpha_ns,\qquad C_{n-1}t=\beta_nt,
\]
where
\[
 \alpha_n=2-2\cos\frac\pi n,\qquad
 \beta_n=2+2\cos\frac\pi n.
\]
Moreover $s_i^2=t_i^2$.  Apply the Rayleigh quotient of $A_w$ to
$x=W^{-1/2}s$ and $y=W^{-1/2}t$.  If
\[
 S_w=s^{\mathsf T}W^{-1}s=t^{\mathsf T}W^{-1}t,
 \qquad r_w=\frac{\alpha_n\|s\|^2}{S_w},
\]
then
\[
 \mathcal Q_{A_w}(x)=r_w,\qquad
 \mathcal Q_{A_w}(y)=\frac{\beta_n}{\alpha_n}r_w.
\]
Thus
$\lambda_{\min}(A_w)\le r_w$ and
$\lambda_{\max}(A_w)\ge(\beta_n/\alpha_n)r_w$, proving
\eqref{eq:path-ratio}.

If equality holds, both Rayleigh inequalities are equalities.  In
particular, $W^{-1/2}s$ is a smallest-eigenvalue eigenvector of $A_w$, so
for some $r>0$,
\[
 W^{1/2}C_{n-1}s=rW^{-1/2}s.
\]
Since $C_{n-1}s=\alpha_ns$ and every $s_i>0$, this says
$\alpha_nw_i=r$ for every $i$.  Conversely, if $W=cI$, then
$A_w=cC_{n-1}$ and equality follows from the displayed endpoint
eigenvalues.
\end{proof}

\begin{remark}[Classical antecedent]\label{rem:forsythe-straus}
The minimization assertion in Proposition~\ref{prop:path-ratio} is not new
as a matrix-scaling theorem.  The matrix $C_{n-1}/2$ has unit diagonal and
Young's property A after the odd and even edge indices are separated.
It is therefore a direct instance of the Forsythe--Straus theorem that
Jacobi scaling is optimal, among positive diagonal congruences, for a
positive definite matrix with Young's property A
\cite[Theorem~4]{ForsytheStraus1955}.  The proof above is included because
the paired sine eigenvectors expose the graph constant and equality case in
a few lines.  Related path weight optimization appears in
\cite{BoydDiaconisSunXiao2006}.
\end{remark}

\section{A spectral-volume consequence}\label{sec:volume}

Put $d=n-1$ and let
\[
 \Delta_{d-1}=\left\{x\in\R_{\ge0}^d:\sum_{i=1}^d x_i=1\right\}.
\]
Following \cite[Section~8]{EJT2026}, the normalized spectral volume $P(G)$
of a connected $n$-vertex graph is the relative Euclidean volume in
$\Delta_{d-1}$ of the closure of its trace-one realizable nonzero spectra,
with all coordinate permutations included.  Strict and nonnegative support
conventions have the same closed spectral region.  Indeed, after trace
normalization the nonnegative weight simplex is compact, its strict interior
is dense, and the ordered spectrum depends continuously on the weights.
We write $B(a,b)=\Gamma(a)\Gamma(b)/\Gamma(a+b)$ for Euler's beta function.

\begin{lemma}[A simplex ratio integral]\label{lem:simplex-ratio}
Let $d\ge2$ and $0<q<1$.  If $X$ is uniform on $\Delta_{d-1}$, then
\begin{equation}\label{eq:simplex-ratio}
 \Pr\!\left(\frac{\min_iX_i}{\max_iX_i}\le q\right)
 =1-\frac{d}{1-q}
 B\!\left(\frac{1+(d-1)q}{1-q},d\right).
\end{equation}
Equivalently, with
$a=(1+(d-1)q)/(1-q)=1+dq/(1-q)$, the right side is
$1-(d-1)B(a,d-1)$.
\end{lemma}

\begin{proof}
Let $E_1,\ldots,E_d$ be independent mean-one exponential random variables.
Then $X_i=E_i/(E_1+\cdots+E_d)$ is uniform on the simplex, and ratios are
unchanged by normalization.  The maximum is unique almost surely.  Choosing
its index and writing its value as $t$ gives
\begin{align*}
 \Pr(\min_iE_i>q\max_iE_i)
 &=d\int_0^\infty e^{-t}
   \left(\int_{qt}^{t}e^{-u}\,du\right)^{d-1}dt\\
 &=d\int_0^\infty e^{-[1+q(d-1)]t}
   (1-e^{-(1-q)t})^{d-1}dt.
\end{align*}
The substitution $z=e^{-(1-q)t}$ gives the beta integral complementary to
\eqref{eq:simplex-ratio}.  The equivalent form follows from
$(a+d-1)B(a,d)=(d-1)B(a,d-1)$.
\end{proof}

\begin{corollary}[Beta-function path bound]\label{cor:path-volume}
For $n\ge3$, set
\[
 q_n=\tan^2\!\left(\frac{\pi}{2n}\right),\qquad
 a_n=\frac{1+(n-2)q_n}{1-q_n}
     =1+\frac{(n-1)q_n}{1-q_n}.
\]
Then
\begin{equation}\label{eq:path-volume-bound}
 \boxed{\quad P(P_n)\le U_n:=1-(n-2)B(a_n,n-2).\quad}
\end{equation}
Equivalently,
\[
 U_n=1-\frac{n-1}{1-q_n}B(a_n,n-1).
\]
Moreover,
\begin{equation}\label{eq:path-volume-asymptotic}
 U_n=\frac{\pi^2}{4n}H_{n-2}
 +O\!\left(\frac{(\log n)^2}{n^2}\right)
 =\frac{\pi^2}{4n}(\log n+\gamma_{\rm E})
 +O\!\left(\frac{(\log n)^2}{n^2}\right).
\end{equation}
In particular, $P(P_n)=O((\log n)/n)$.
\end{corollary}

\begin{proof}
A realizable normalized positive spectrum of $P_n$ has $d=n-1$ coordinates
and, by Proposition~\ref{prop:path-ratio}, satisfies
\[
 \frac{\min_i\lambda_i}{\max_i\lambda_i}\le q_n.
\]
Lemma~\ref{lem:simplex-ratio} gives the exact volume of this necessary
region and hence \eqref{eq:path-volume-bound}.

For the asymptotic write $a_n=1+\epsilon_n$.  Since
\[
 q_n=\frac{\pi^2}{4n^2}+O(n^{-4}),\qquad
 \epsilon_n=\frac{(n-1)q_n}{1-q_n}
 =\frac{\pi^2}{4n}+O(n^{-2}),
\]
we have
\[
 1-U_n=(n-2)B(1+\epsilon_n,n-2)
 =\frac{\Gamma(1+\epsilon_n)\Gamma(n-1)}
        {\Gamma(n-1+\epsilon_n)}.
\]
Using $\psi(n-1)+\gamma_{\rm E}=H_{n-2}$ and Taylor expansion of
$\log\Gamma$ gives
\[
 \log(1-U_n)=-\epsilon_nH_{n-2}+O(\epsilon_n^2).
\]
Since $\epsilon_nH_{n-2}=O((\log n)/n)$, exponentiation proves
\eqref{eq:path-volume-asymptotic}.
\end{proof}

\begin{remark}[Checks and scope]
For $n=2$, $P_2=K_2$ realizes every one-point positive spectrum, so
$P(P_2)=1$.  For $n=3$, $q_3=1/3$, $a_3=2$, and $U_3=1/2$; this is exact
by the $P_3$ characterization in \cite{FGL2025}.  For $n=4$ the bound is
$U_4=9/17$, but the exact $P_4$ region has further constraints
\cite{EJT2026}.  The asymptotic equivalence in
\eqref{eq:path-volume-asymptotic} belongs to the upper bound $U_n$, not to
$P(P_n)$ itself.
\end{remark}

\section{Exact lifting to hyperbolic surfaces}\label{sec:surface}

We now transfer regular graph charts to closed hyperbolic surfaces of fixed
genus.  Throughout this section the surface eigenvalues are indexed by
\[
 0=\lambda_0(X)<\lambda_1(X)\le\lambda_2(X)\le\cdots.
\]

Fix $g\ge2$ and put
\begin{equation}\label{eq:N-m-genus}
 N=2g-2,\qquad m=N-1=2g-3.
\end{equation}
Let $T_g$ be a tree with $g$ vertices of degree one and $g-2$ vertices of
degree three.  Thus $|V(T_g)|=2g-2$ and $|E(T_g)|=2g-3$; for $g=2$ it is a
single edge.

Associate a one-holed torus to each leaf and a pair of pants to each
trivalent vertex, then glue boundary components according to the edges of
$T_g$.  The number of available boundary components is
\[
 g+3(g-2)=4g-6=2|E(T_g)|,
\]
so all are paired.  Every piece has Euler characteristic $-1$, and hence the
closed connected result has
\[
 \chi=-g-(g-2)=2-2g.
\]
It is a surface of genus $g$.

This decomposition is adapted to Burger's weighted graph.  If $Y_v$ is a
piece, its graph vertex mass in the normalization of \cite{Burger1990} is
\[
 \mathfrak m(v)=-\chi(Y_v)=1
\]
for both a pair of pants and a one-holed torus.  Equivalently, every piece
has area $2\pi$.  The limiting operator is therefore the ordinary weighted
combinatorial Laplacian with unit vertex masses.

For $w\in\R_{>0}^{m}$ let
\[
 L_{T_g}(w)=\sum_{e=uv\in E(T_g)}
 w_e(\mathbf e_u-\mathbf e_v)(\mathbf e_u-\mathbf e_v)^{\mathsf T},
\]
and write its positive eigenvalues as
\[
 0<\mu_1(w)\le\cdots\le\mu_m(w),\qquad
 \Lambda_{T_g}(w)=(\mu_1(w),\ldots,\mu_m(w)).
\]

Choose the remaining Fenchel--Nielsen data once and for all.  Inside every
one-holed torus, fix the length and twist of a nonseparating curve, and fix
the twists along the separating curves represented by $T_g$.  For $t>0$
and $w>0$, let $X(t,w)$ be the closed curvature-$-1$ surface whose
separating cuff associated to $e$ has length
\begin{equation}\label{eq:cuff-length}
 \ell_e(t,w)=tw_e.
\end{equation}
These are all $3g-3$ length coordinates and all $3g-3$ twist coordinates of
a pants decomposition.  Thus $X(t,w)$ is a well-defined marked hyperbolic
surface, depending real-analytically on $(t,w)$ for $t>0$.

Let $A(t,w)$ be the multicurve formed by the $m$ separating cuffs.  We first
record the uniform geometric input needed in Burger's comparison theorem.
Following \cite{Burger1990}, write
\begin{equation}\label{eq:burger-L-definition}
 L(A)=\min\!\left\{2\operatorname{arsinh}(1),\ 
  \ell(\gamma):\gamma\text{ is a closed geodesic disjoint from }A
  \text{ and not peripheral to }A\right\},
\end{equation}
with the first entry retained when the second set is empty, and put
$l(A)=\max_{\alpha\in A}\ell(\alpha)$.

\begin{lemma}[Uniform nonpinching scale]\label{lem:nonpinching-scale}
For every compact $C\Subset\R_{>0}^{m}$ there are $t_C,c_C>0$ such that,
for $0<t<t_C$ and $w\in C$, Burger's nonpinching length scale satisfies
\[
  L(A(t,w))\ge c_C.
\]
Consequently Burger's cusp parameter obeys
\begin{equation}\label{eq:burger-cusp-parameter}
  \epsilon(t,w):=\frac{l(A(t,w))}{L(A(t,w))}
  \le \frac{t\max_{w\in C,e}w_e}{c_C}=O_C(t).
\end{equation}
\end{lemma}

\begin{proof}
The assertion is immediate on every parameter subinterval bounded away from
$t=0$, by compactness of Fenchel--Nielsen coordinates.  Consider therefore a
sequence with $t\downarrow0$.  A closed geodesic disjoint from $A(t,w)$ lies
in one complementary piece.  A pair of pants contains no nonperipheral
essential simple closed curve.  Any closed geodesic shorter than
$2\operatorname{arsinh}(1)$ is simple, so a sequence of closed geodesics
with length tending to zero is eventually simple.  A one-holed torus has
fixed positive internal length and fixed twist coordinates, while its
boundary length tends to zero in a bounded interval.  The fixed positive
internal length excludes internal nodal strata.  These tori therefore remain
in a compact subset of the augmented one-holed-torus moduli space, and their
nonperipheral systoles have a uniform positive lower bound.  A curve meeting
a component of $A(t,w)$ cannot become short, by the collar lemma.  These
observations give the required lower bound for the nonpinching scale; see,
for example,
\cite[Chapters~3--4]{Buser1992}.  Since
$l(A(t,w))=t\max_e w_e$, \eqref{eq:burger-cusp-parameter} follows.
\end{proof}

\begin{lemma}[Uniform Burger limit]\label{lem:burger-uniform}
For every compact $C\Subset\R_{>0}^{m}$,
\begin{equation}\label{eq:burger-uniform}
 \lim_{t\downarrow0}\sup_{w\in C}\max_{1\le j\le m}
 \left|\frac{2\pi^2}{t}\lambda_j(X(t,w))-\mu_j(w)\right|=0.
\end{equation}
\end{lemma}

\begin{proof}
For the partition by the $m$ separating cuffs, Burger's graph is $T_g$,
every vertex mass is one, and the conductance vector is $tw$.  Its positive
eigenvalues, after discarding the common zero mode, are
$t\mu_1(w),\ldots,t\mu_m(w)$.  Lemma~\ref{lem:nonpinching-scale} gives
$\epsilon=O_C(t)$ uniformly.  Hence, for $t$ small enough that Burger's
condition $\epsilon<\alpha_1$ holds throughout $C$, the uniform comparison
\cite[Theorem~1.1]{Burger1990} gives, in the standard combinatorial
normalization,
\[
 \frac{\lambda_j(X(t,w))}{t\mu_j(w)}
 =\frac1{2\pi^2}(1+o(1)).
\]
The error is uniform in a fixed topological cusp neighborhood; it is
controlled there by a constant times
$\sqrt\epsilon+\epsilon|\log\epsilon|$.  The functions $\mu_j(w)$ are
bounded on $C$, so the multiplicative estimate implies the absolute
convergence \eqref{eq:burger-uniform}.  In addition, Otal--Rosas gives the
pointwise bound $\lambda_{m+1}=\lambda_{2g-2}>1/4$, identifying these as
exactly the collapsing positive modes.
\end{proof}

\begin{remark}[The normalization constant]\label{rem:normalization}
The factor $2\pi^2$ must be retained when $L_{T_g}(w)$ denotes the standard
combinatorial Laplacian.  Burger's vertex mass is one here, whereas the
hyperbolic area of each piece is $2\pi$.  Equivalently, the graph quadratic
form $\pi^{-1}\sum_e w_e(F_u-F_v)^2$ with scalar product
$2\pi\sum_vF_vG_v$ has operator $(2\pi^2)^{-1}L_{T_g}(w)$.  Mixing these two
conventions would change the prescribed scale.
\end{remark}

The next theorem is the exact-lifting step.  It uses only continuous
dependence of ordered eigenvalues and Brouwer degree, not differentiability
of the surface spectral map at the nodal limit.

\begin{theorem}[Regular graph charts give exact surface spectra]
\label{thm:surface-transfer}
Suppose $\rho=(\rho_1,\ldots,\rho_m)$, with
$0<\rho_1<\cdots<\rho_m$, has a strict regular realization on $T_g$:
there is $w^0>0$ such that
\[
 \Lambda_{T_g}(w^0)=\rho,
 \qquad \det D\Lambda_{T_g}(w^0)\ne0.
\]
Then there are a neighborhood $V_\rho$ of $\rho$ and $s_\rho>0$ such that,
for every $\eta\in V_\rho$ and every $0<s<s_\rho$, some closed hyperbolic
surface $X$ of genus $g$ satisfies
\begin{equation}\label{eq:exact-surface}
 (\lambda_1(X),\ldots,\lambda_m(X))=s\eta.
\end{equation}
The threshold may be chosen so that $s\eta_m<1/4$ throughout this set.
\end{theorem}

\begin{proof}
The inverse-function theorem gives a connected neighborhood $U_0$ of $w^0$
on which $\Lambda_{T_g}$ is a diffeomorphism onto its image.  Shrink $U_0$
inside the simple-spectrum locus, and choose $r>0$ so that
\[
 \overline{B_{2r}(\rho)}\subset
 \Lambda_{T_g}(U_0)\cap\{0<\eta_1<\cdots<\eta_m\}.
\]
Set
\[
 D=(\Lambda_{T_g}|_{U_0})^{-1}(B_{2r}(\rho)).
\]
After shrinking $U_0$, we have
$\overline D\Subset\R_{>0}^{m}$, and $\Lambda_{T_g}$ maps $D$
diffeomorphically onto $B_{2r}(\rho)$.

For $t>0$ define the continuous map
\[
 \Phi_t(w)=\frac{2\pi^2}{t}
 (\lambda_1(X(t,w)),\ldots,\lambda_m(X(t,w))),
 \qquad w\in\overline D.
\]
Continuity follows from the Fenchel--Nielsen construction of $X(t,w)$ and
the min--max characterization of ordered Laplace eigenvalues.
Lemma~\ref{lem:burger-uniform} gives
\begin{equation}\label{eq:uniform-chart-convergence}
 \sup_{w\in\overline D}\|\Phi_t(w)-\Lambda_{T_g}(w)\|\longrightarrow0.
\end{equation}
Put $V_\rho=B_r(\rho)$.  For $\eta\in V_\rho$ and $w\in\partial D$,
$\|\Lambda_{T_g}(w)-\eta\|\ge r$.  Consequently, for all sufficiently
small $t$, uniformly in $\eta\in V_\rho$, the straight-line homotopy from
$\Lambda_{T_g}$ to $\Phi_t$ avoids $\eta$ on $\partial D$.  Homotopy
invariance of Brouwer degree gives
\[
 \deg(\Phi_t,D,\eta)
 =\deg(\Lambda_{T_g},D,\eta)
 =\operatorname{sgn}\det D\Lambda_{T_g}(w^0)\ne0.
\]
For the last equality, the graph map has one preimage of $\eta$ in $D$, and
the sign of its nonzero Jacobian is constant on the connected set $D$.
Thus $\Phi_t(w(t,\eta))=\eta$ for some $w(t,\eta)\in D$, or
\[
 \lambda_j(X(t,w(t,\eta)))=\frac{t}{2\pi^2}\eta_j.
\]
Taking $t=2\pi^2s$ proves \eqref{eq:exact-surface}; shrink $s_\rho$ to
obtain the final inequality.
\end{proof}

The graph base point may depend on $\rho$ without causing a uniformity
problem.  For a fixed target it is chosen once, and the same degree domain
works at every sufficiently small scale; the realizing weight is allowed to
vary with the scale.  The argument supplies existence at each scale, not a
continuous selection of realizing weights or surfaces.

\begin{corollary}[Compact-uniform exact lift]\label{cor:compact-lift}
Let $\Omega(T_g)$ be the set of simple spectra admitting a strict regular
realization on $T_g$.  For every compact $K\Subset\Omega(T_g)$ there is
$s_{g,K}>0$ such that, for every $\rho\in K$ and $0<s<s_{g,K}$, some closed
hyperbolic surface $X$ of genus $g$ satisfies
\[
 (\lambda_1(X),\ldots,\lambda_{2g-3}(X))=s\rho.
\]
\end{corollary}

\begin{proof}
The neighborhoods $V_\rho$ from Theorem~\ref{thm:surface-transfer} cover
$K$.  Take a finite subcover and the minimum of the associated positive
thresholds.
\end{proof}

\section{Full-index prescription in genus \texorpdfstring{$g$}{g}}
\label{sec:full-index}

\begin{proof}[Proof of Theorem~\ref{thm:intro-surface}]
Take
\[
 \varepsilon_g=\delta_{2g-2}
\]
from Theorem~\ref{thm:intro-universal}.  The tree $T_g$ is connected on
$2g-2$ vertices, so every $\rho\in\cL_{2g-2}(\varepsilon_g)$ has a strict
regular realization on $T_g$.  Theorem~\ref{thm:surface-transfer} gives the
exact equality for every sufficiently small $s$.  The theorem of
Otal--Rosas \cite{OtalRosas2009} states that
\begin{equation}\label{eq:otal-rosas}
 \lambda_{2g-2}(X)>\frac14
\end{equation}
for every closed hyperbolic surface of genus $g$.  After reducing the
threshold so that $s\rho_{2g-3}<1/4$, the prescribed list is therefore
precisely the full collection of eigenvalues below $1/4$, and the next
eigenvalue satisfies the asserted strict $1/4$ bound.
\end{proof}

\begin{corollary}[A compact lacunary region]\label{cor:lacunary-compact}
If $K\Subset\cL_{2g-2}(\varepsilon_g)$, there is one $s_{g,K}>0$ such that
every $s\rho$, with $\rho\in K$ and $0<s<s_{g,K}$, is exactly the first
$2g-3$ positive-eigenvalue list of a closed genus-$g$ hyperbolic surface.
\end{corollary}

\begin{proof}
Combine Theorem~\ref{thm:intro-universal} with
Corollary~\ref{cor:compact-lift}, reducing the threshold using
\eqref{eq:otal-rosas}.
\end{proof}

This gives an exact general-genus near-origin region for
\cite[Question~1]{EJT2026}.  This lacunary corollary is a partial answer: it
neither characterizes the entire cube $[0,1/4]^{2g-3}$ nor treats lists with
repeated entries.  The
exact graph-to-surface degree mechanism is classical: it is the
essential-family construction of Colbois--Colin de Verdi\`ere
\cite[Lemma~V.1 and Section~VI]{ColboisCdV1988}, used there with a complete
graph to create high multiplicity.  The ingredient specific to the present
paper is the universal regular lacunary chart.  Choosing it on the
genus-compatible tree $T_g$ reaches the full admissible length $2g-3$.

\section{The complexity-one stable-graph spectral link}
\label{sec:tangent-link}

We now compactify the preceding single graph chart.  A
\emph{complexity-one stable genus-$g$ pseudograph} is a connected
pseudograph $\Gamma$, with loops and parallel edges allowed, having
$N=2g-2$ vertices, arithmetic genus
\begin{equation}\label{eq:stable-arithmetic-genus}
 b_1(\Gamma)+\sum_v h(v)=g,
\end{equation}
and a genus label $h(v)\ge0$ at every vertex such that
\begin{equation}\label{eq:complexity-one-vertex}
 2h(v)-2+\deg_\Gamma(v)=1.
\end{equation}
Loops count twice in the degree.  Thus a vertex represents either a pair of
pants, $(h,\deg)=(0,3)$, or a one-holed torus, $(h,\deg)=(1,1)$.  There are
only finitely many isomorphism types; denote their collection by
$\mathfrak G_g$.

Loops remain part of the topology but contribute no energy to functions
constant on a component.  Combine the conductances of all parallel nonloop
cuffs joining the same two vertices, and call the resulting set of distinct
pairs $E_{\rm eff}(\Gamma)$.  On the compact trace-one simplex put
\[
 K_\Gamma=\left\{w\in\R_{\ge0}^{E_{\rm eff}(\Gamma)}:
                 2\sum_e w_e=1\right\}.
\]
If $0=\mu_1\le\cdots\le\mu_N$ are the eigenvalues of $L_\Gamma(w)$,
including any extra zeros caused by the positive support disconnecting, set
\[
 \psi_\Gamma(w)=(\mu_2,\ldots,\mu_N).
\]

We shall use the comparison theorem also when the cuff lengths have
arbitrarily different orders.  The following formulation isolates the
needed uniformity.

\begin{lemma}[Uniform projectivized Burger comparison]
\label{lem:burger-projective}
Let $X_k$ converge in the augmented moduli space after pinching a nonempty
full Deligne--Mumford multicurve $A_k$, and suppose that the complement has
$r$ components.  Let $G_k$ be Burger's weighted graph, with vertex masses
$-\chi(Y_{v,k})$ and cuff-length conductances.  Then, after passing to the
convergent augmented chart,
\[
 L(A_k)\ge c>0,\qquad \epsilon_k=\frac{l(A_k)}{L(A_k)}\longrightarrow0,
\]
and
\begin{equation}\label{eq:burger-next-mode}
 \liminf_{k\to\infty}\lambda_r(X_k)>0.
\end{equation}
If $r\ge2$ and
$0<\nu_{1,k}\le\cdots\le\nu_{r-1,k}$ are the positive graph eigenvalues,
then
\begin{equation}\label{eq:burger-projective}
 \max_{1\le j\le r-1}
 \left|\frac{2\pi^2\lambda_j(X_k)}{\nu_{j,k}}-1\right|
 \longrightarrow0.
\end{equation}
Thus exactly $r-1$ positive surface eigenvalues tend to zero.  The estimate
is independent of the ratios among the pinching cuff lengths.  Consequently,
when $r\ge2$ and the graph spectral trace is normalized, the projectivized
surface and graph spectra have the same subsequential limits, including
limits on weight-simplex faces.
\end{lemma}

\begin{proof}
Fullness means that, on each component of the limiting noded surface, no
additional essential geodesic disjoint from $A_k$ tends to zero.  Compactness
in the chosen augmented chart and the definition
\eqref{eq:burger-L-definition} therefore give $L(A_k)\ge c>0$, whereas
$l(A_k)\to0$.  Burger's multiplicative estimate
\cite[Theorem~1.1]{Burger1990} depends only on the topology and
$\epsilon_k$; it does not involve the smallest-to-largest conductance ratio.
It applies simultaneously to the $r-1$ positive graph modes and gives
\eqref{eq:burger-projective} when $r\ge2$.  Because every cuff in $A_k$
shrinks, all these graph eigenvalues tend to zero.  On the orthogonal
complement of the componentwise constants, the limiting components have a
positive Poincar\'e threshold; compactness of the augmented chart makes that
threshold uniform.  The min--max principle, equivalently the noncollapsing
part of Burger's comparison, gives \eqref{eq:burger-next-mode}.  Hence there
are exactly $r-1$ collapsing positive modes.  Dividing the collapsing
spectral lists by their sums cancels the common constant and proves the last
assertion.
\end{proof}

\begin{theorem}[Exact maximal-cusp tangent link]
\label{thm:tangent-link}
For every $g\ge2$,
\begin{equation}\label{eq:tangent-link}
 \boxed{\quad
 \mathcal T_g=\bigcup_{\Gamma\in\mathfrak G_g}
                    \psi_\Gamma(K_\Gamma).
 \quad}
\end{equation}
In particular, $\mathcal T_g$ is compact and semialgebraic.  Equivalently,
without fixing trace, the right side consists of
\[
 \frac{(\mu_2(L_\Gamma(w)),\ldots,\mu_N(L_\Gamma(w)))}
      {\tr L_\Gamma(w)},
 \qquad w\ge0,\quad \tr L_\Gamma(w)>0.
\]
The zero conductances in \eqref{eq:tangent-link} record multiscale pinching;
they cannot in general be discarded.
\end{theorem}

\begin{proof}
Suppose $\lambda_m(X_k)\to0$.  The sequence leaves every compact part of
moduli space, so after passing to a Deligne--Mumford subsequence it has a
nonempty full pinching multicurve $A_k$.  If its
complement has $r$ components, Burger's cusp-partition theorem and
Lemma~\ref{lem:burger-projective} give precisely $r-1$ collapsing graph
modes.  Thus $r-1\ge m$.  On the other hand, stability gives
\[
 \sum_{v=1}^r-\chi(Y_v)=2g-2=N,
\]
and every summand is a positive integer.  Hence $r\le N=m+1$, so equality
holds throughout: $r=N$, there are exactly $m$ collapsing positive modes,
and every component has mass one.  Each component is therefore a pair of
pants or a one-holed torus, and the stable dual type belongs to
$\mathfrak G_g$.  Otal--Rosas supplies independently
$\lambda_{m+1}(X_k)>1/4$.

Pass to a subsequence on which the finite graph type is fixed.  After
aggregating parallel cuffs and normalizing total conductance, compactness of
$K_\Gamma$ gives $w_k\to w$.  Burger's multiplicative comparison, with the
common normalization cancelled, remains valid as the normalized weights
approach a face by Lemma~\ref{lem:burger-projective}.  It yields
$p(X_k)\to\psi_\Gamma(w)$.  Vanishing coordinates of $w$ retain the extra
zero graph eigenvalues and therefore cover multiscale limits.  This proves
one inclusion.

Conversely, fix $(\Gamma,w)$ on the right.  Approximate every zero
conductance by trace-one vectors $w^{(a)}>0$ with $w^{(a)}\to w$, split an
effective conductance
positively among parallel cuffs, and choose positive lengths for loop cuffs.
Realize the vertices by their complexity-one hyperbolic pieces, freeze the
remaining Fenchel--Nielsen data, and multiply all cuff lengths by $t$.
For each $a$, Lemma~\ref{lem:burger-projective} gives the direction
$\psi_\Gamma(w^{(a)})$ as $t\downarrow0$.  Choose $t_a>0$ so that the
projectivized error is below $1/a$ and $\lambda_m<1/a$.  Then the diagonal
sequence converges to $\psi_\Gamma(w)$ while $\lambda_m\to0$.

Finally, compactness follows from the finite union in
\eqref{eq:tangent-link}.  Semialgebraicity follows by expressing the ordered
eigenvalues as ordered roots of the characteristic polynomial and applying
Tarski--Seidenberg to each compact weight simplex.
\end{proof}

\begin{corollary}[Off-link exclusion]\label{cor:off-link}
If $p=(p_1,\ldots,p_m)$ satisfies
$0<p_1\le\cdots\le p_m$ and $\sum_jp_j=1$, and
$p\notin\mathcal T_g$, then there is
$s_p>0$ such that $sp$ does not occur as the first $m$ positive-eigenvalue
list of a genus-$g$ closed hyperbolic surface for any $0<s<s_p$.
\end{corollary}

\begin{proof}
Otherwise choose $s_k\downarrow0$ and surfaces whose first $m$
positive-eigenvalue lists are $s_kp$.
Their normalized directions equal $p$ and their $m$th eigenvalues tend to
zero, contradicting the definition of $\mathcal T_g$.
\end{proof}

\begin{corollary}[Compact-uniform off-link exclusion]
\label{cor:off-link-compact}
Let $K$ be a compact subset of
\[
 \left\{0<p_1\le\cdots\le p_m:\sum_jp_j=1\right\}
 \setminus\mathcal T_g.
\]
There is one $s_K>0$ such that no $sp$, with $p\in K$ and
$0<s<s_K$, occurs as the first $m$ positive-eigenvalue list of a closed
genus-$g$ hyperbolic surface.
\end{corollary}

\begin{proof}
Otherwise there are $p_k\in K$, $s_k\downarrow0$, and realizing surfaces.
After taking a subsequence, $p_k\to p\in K$.  Their $m$th eigenvalues tend
to zero, so the definition of the tangent link gives $p\in\mathcal T_g$, a
contradiction.
\end{proof}

\section{Generic exactness and the singular skeleton}
\label{sec:generic}

We allow every relative-open face of every $K_\Gamma$.  A point $w$ on such
a face is \emph{normalized regular} if $\psi_\Gamma(w)$ has positive,
pairwise distinct coordinates and the derivative along the face has rank
$m-1$.  Define
\begin{equation}\label{eq:U-B-def}
 \mathcal U_g=\{p\in\mathcal T_g:p\text{ has a normalized-regular
 preimage on some face}\},\qquad
 \mathfrak B_g=\mathcal T_g\setminus\mathcal U_g.
\end{equation}
This definition is deliberately existential.  A spectral value can have a
critical preimage and a regular preimage simultaneously, so the image of all
critical weights is not the boundary and is not the correct skeleton.

We first extend Theorem~\ref{thm:surface-transfer} from the chosen tree to a
stable graph.  Split parallel effective weights in fixed positive
proportions and freeze all loop and unused cuff variables.

\begin{lemma}[Uniform convergence in a stable chart]
\label{lem:stable-chart-uniform}
Fix $\Gamma\in\mathfrak G_g$, a positive splitting of each effective
conductance among its parallel cuffs, positive rescaled lengths for all loop
cuffs, and compact thick and twist data on the complementary pieces.  If
$C\Subset\R_{>0}^{E_{\rm eff}(\Gamma)}$, let $X_\Gamma(t,w)$ be the surface
obtained by multiplying all nodal cuff lengths by $t$.  Writing
$\Lambda_\Gamma(w)=(\mu_2(w),\ldots,\mu_N(w))$, one has
\begin{equation}\label{eq:stable-chart-uniform}
 \sup_{w\in C}\left\|
 \frac{2\pi^2}{t}\Lambda_g(X_\Gamma(t,w))-\Lambda_\Gamma(w)
 \right\|\longrightarrow0
 \qquad(t\downarrow0).
\end{equation}
The same conclusion is uniform when the fixed auxiliary data range over a
compact subset of the positive parameter locus; in particular, all splitting
proportions and rescaled loop-cuff lengths then stay uniformly away from zero.
\end{lemma}

\begin{proof}
All complementary components have complexity one and hence Burger mass one.
The chosen positive splittings sum to the prescribed effective
conductances, whereas loop cuffs contribute no graph energy.  Compactness of
$C$ and of the auxiliary data gives a uniform lower bound for $L(A)$ by the
argument of Lemma~\ref{lem:nonpinching-scale}; the cusp parameter is therefore
$O(t)$.  Burger's topology-uniform multiplicative comparison then gives
\eqref{eq:stable-chart-uniform}.  The final assertion follows from the same
estimates on the product compact set.
\end{proof}

\begin{proposition}[Stable-graph regular transfer]
\label{prop:stable-transfer}
Let $\Gamma\in\mathfrak G_g$ have a strict effective weighting $w^0$ with
simple positive spectrum $\rho$.  If the unnormalized spectral derivative
has rank $m$, then there are a neighborhood $V_\rho$ and $s_\rho>0$ such
that every $s\eta$, $\eta\in V_\rho$ and $0<s<s_\rho$, is exactly the first
$m$ positive-eigenvalue list of a genus-$g$ closed hyperbolic surface.
\end{proposition}

\begin{proof}
Choose an invertible $m\times m$ spectral minor and vary only those cuff
weights on a small affine slice, keeping all other positive weights fixed.
On a compact subset of that slice,
Lemma~\ref{lem:stable-chart-uniform} makes the scaled surface eigenvalue map
converge uniformly to the graph map.  The latter has local Brouwer degree
$\pm1$.  Choosing a degree domain and a smaller target ball exactly as in
Theorem~\ref{thm:surface-transfer}, uniform convergence makes the straight
homotopy boundary-avoiding.  Nonzero degree gives the desired preimage.
Otal--Rosas separates the next surface eigenvalue from the collapsing
cluster.
\end{proof}

\begin{theorem}[Generic exact maximal-cusp lifting]
\label{thm:generic-exact}
For every $g\ge2$, $\mathcal U_g$ is relatively open and semialgebraic in the
strict ordered simplex (rather than merely relatively open in
$\mathcal T_g$), and $\mathfrak B_g$ is compact semialgebraic with
\begin{equation}\label{eq:skeleton-dimension}
 \dim\mathfrak B_g\le m-2,\qquad
 \mathcal H^{m-1}(\mathfrak B_g)=0.
\end{equation}
Here $\dim\varnothing=-\infty$; in particular,
$\mathfrak B_2=\varnothing$.
For every $p\in\mathcal U_g$ there are a relative neighborhood $V_p$ and
$s_p>0$ such that every $s\eta$, $\eta\in V_p$ and $0<s<s_p$, is exactly
realized by a genus-$g$ closed hyperbolic surface.  Consequently almost every
direction in $\mathcal T_g$, with respect to ambient
$\mathcal H^{m-1}$ measure in the trace-one hyperplane, has an exact
all-small-scale ray lift.
\end{theorem}

\begin{proof}
For each relative-open face $F$ of $K_\Gamma$, restrict the normalized map to
its simple-spectrum locus and denote it by $\psi_{\Gamma,F}$.  Let
$\Crit\psi_{\Gamma,F}$ be the set where its derivative has rank less than
$m-1$; this is the whole source when $\dim F<m-1$.  Let
$\Delta_{\rm sing}$ be the union of the coordinate and collision walls in
the ordered trace-one simplex.  The existential definition of regularity
gives the explicit containment
\begin{equation}\label{eq:skeleton-critical-containment}
 \mathfrak B_g\subseteq
 \Delta_{\rm sing}\ \cup\!
 \bigcup_{\Gamma\in\mathfrak G_g}\ \bigcup_{F\preceq K_\Gamma}
 \psi_{\Gamma,F}(\Crit\psi_{\Gamma,F}).
\end{equation}
Semialgebraic Sard theory \cite{KurdykaOrroSimon2000} puts every critical
value set on the right in dimension at most $m-2$; faces of smaller
dimension satisfy the same bound directly.  The walls also have dimension
at most $m-2$.  Since the union is finite,
\eqref{eq:skeleton-dimension} follows.

Characteristic coefficients, Jacobian minors, and quantifier elimination
show that $\mathcal T_g$, $\mathcal U_g$, and $\mathfrak B_g$ are
semialgebraic.  At a normalized submersion, the submersion theorem supplies
a neighborhood in the ambient trace-one hyperplane, all of whose points
retain regular preimages.  Hence $\mathcal U_g$ is open in the strict ordered
simplex.  Its complement in the compact set $\mathcal T_g$ is therefore
closed, proving compactness of $\mathfrak B_g$.  For $g=2$, both the tangent
link and the trace-one target are the singleton $(1)$, whose zero-dimensional
map is regular, so the skeleton is empty.  Finally, the lacunary tree chart
gives a full-dimensional subset of $\mathcal T_g$; thus the ambient-measure
form of ``almost every'' is nonvacuous.

Fix a regular preimage.  Homogeneity adds the radial direction and turns
normalized rank $m-1$ into unnormalized rank $m$.  Because $p_1>0$,
$\ker L_\Gamma(w)=\R\one$; equivalently, the positive effective support is a
connected spanning subgraph.  If some allowed graph weights vanish, the
regular spanning-subgraph lift Lemma~\ref{lem:supergraph} makes every new
effective adjacency positive while preserving the spectrum and rank.
Parallel weights can then be split and loops assigned arbitrary positive
lengths.  Proposition~
\ref{prop:stable-transfer} gives the claimed exact neighborhood and ray.
\end{proof}

\begin{corollary}[Compact-uniform wedges]\label{cor:generic-compact}
If $K\Subset\mathcal U_g$, then one $s_K>0$ works for every
$p\in K$ and $0<s<s_K$.
\end{corollary}

\begin{proof}
Take a finite subcover of $K$ by the neighborhoods in
Theorem~\ref{thm:generic-exact} and the minimum of their scale thresholds.
\end{proof}

The following elementary inclusion is useful when reading the skeleton.
If $F:K\to\R^d$ is $C^1$ on a compact manifold with corners, and
$\Crit F$ denotes the interior points where $DF$ is not onto, then
\begin{equation}\label{eq:image-boundary-inclusion}
 \partial F(K)\subseteq F(\partial K\cup\Crit F).
\end{equation}
Indeed, an interior preimage at which $DF$ is onto makes the image point
interior by the submersion theorem.  The converse to
\eqref{eq:image-boundary-inclusion} is false: the existence of a boundary or
critical preimage does not make its value a boundary point.

\section{Uniformly conditioned clusters and finite-prefix nonrigidity}
\label{sec:conditioned-prefixes}

The regular part of the maximal-cusp link contains a full-dimensional region
whose internal condition number is bounded independently of the genus.  The
input is the finite-size strict form of the bipartite Ramanujan theorem.

\begin{theorem}[An epsilon-free Ramanujan wedge]
\label{thm:ramanujan-wedge}
For every $g\ge2$, with $m=2g-3$, there are a nonempty open set
\[
 V_g\subset\{0<\eta_1<\cdots<\eta_m\}
\]
and $s_g>0$ such that every $s\eta$, $\eta\in V_g$ and $0<s<s_g$, is
exactly the first $m$ positive-eigenvalue list of a closed genus-$g$
hyperbolic surface, with
\begin{equation}\label{eq:ramanujan-condition}
 \frac{\eta_m}{\eta_1}<18+12\sqrt2.
\end{equation}
Otal--Rosas gives $\lambda_{m+1}=\lambda_{2g-2}>1/4$ for every realizing
surface.  The threshold is reduced only to ensure $s\eta_m<1/4$; hence the
prescribed list is the entire spectrum below $1/4$.
\end{theorem}

\begin{proof}
Write $N=2g-2=2d$.  Theorem~1.1 of
Marcus--Spielman--Srivastava \cite{MSS2018}, applied with degree three and
bipartition size $d+d$, gives a cubic bipartite multigraph $\Gamma$ on $N$
vertices whose nontrivial adjacency eigenvalues $\alpha$ satisfy the strict
finite-size estimate
\[
 |\alpha|<2\sqrt2.
\]
The eigenvalue $3$ is therefore simple, so $\Gamma$ is connected.  Moreover
\[
 |E(\Gamma)|=\frac{3N}{2}=3g-3,\qquad
 b_1(\Gamma)=|E|-|V|+1=g.
\]
Thus $\Gamma$ is an all-pants complexity-one stable genus-$g$ pseudograph;
parallel
edges are legitimate cuff edges.

Give every cuff unit conductance and aggregate the multiplicities between
each vertex pair.  The resulting effective weighted Laplacian is
$L=3I-A$.  The trivial bipartite adjacency eigenvalue $-3$ gives
$\mu_{\max}=6$, while
\[
 \mu_{\min}>3-2\sqrt2.
\]
Consequently
\begin{equation}\label{eq:ramanujan-ratio-calculation}
 \frac{\mu_{\max}}{\mu_{\min}}
 <\frac6{3-2\sqrt2}=18+12\sqrt2.
\end{equation}

Let $\widehat\Gamma$ be the connected effective support after parallel
edges are aggregated.  Theorem~\ref{thm:intro-universal} supplies a simple,
coefficient-regular positive weighting on $\widehat\Gamma$.  Hence some
$m$-row coefficient-Jacobian minor and the spectral discriminant are
nonzero polynomials on its weight space.  Their simultaneous nonvanishing
is dense.  We may therefore perturb the aggregated unit-conductance vector
arbitrarily slightly to a positive weighting $w^*$ having simple spectrum
and unnormalized spectral rank $m$.  The strict margin in
\eqref{eq:ramanujan-ratio-calculation} preserves
\eqref{eq:ramanujan-condition}.  Split each effective conductance into
fixed positive conductances on its parallel cuffs and apply
Proposition~\ref{prop:stable-transfer}.  Shrinking the resulting
neighborhood preserves the strict ratio, and Otal--Rosas supplies the last
assertions.  For $g=2$ the graph is the two-vertex triple-edge graph; its
single positive eigenvalue has condition number one and the same argument is
immediate.
\end{proof}

\begin{corollary}[A maximal cluster with arbitrary outside gap]
\label{cor:maximal-outside-gap}
For every $g\ge2$ and every $M>0$, there is a closed genus-$g$ hyperbolic
surface such that
\[
 \frac{\lambda_{2g-3}}{\lambda_1}<18+12\sqrt2,
 \qquad
 \frac{\lambda_{2g-2}}{\lambda_{2g-3}}>M.
\]
\end{corollary}

\begin{proof}
Choose $\eta\in V_g$ and then
$0<s<\min\{s_g,1/(4M\eta_m)\}$.  Theorem~\ref{thm:ramanujan-wedge} gives
$\lambda_m=s\eta_m$ and $\lambda_{m+1}>1/4$.
\end{proof}

We next make explicit the parameter uniformity behind finite-prefix
nonrigidity.  The statement uses normalized twist angles and rescaled cuff
weights in a local augmented Fenchel--Nielsen chart; these stay in a compact
set as the common pinching scale tends to zero.

\begin{proposition}[Parametric stable-graph transfer]
\label{prop:parametric-transfer}
Let $\Gamma$ be an all-pants complexity-one stable genus-$g$ graph and let
$w^0>0$
be a simple, unnormalized-regular effective weighting with spectrum $\rho$.
There are a neighborhood $V_\rho$, a number $s_\rho>0$, and a compact
parameter cell $Q$ of dimension
\[
 (6g-6)-(2g-3)=4g-3
\]
with nonempty interior such that, for every $\eta\in V_\rho$,
$0<s<s_\rho$, and $q\in Q$, some surface with complementary parameter $q$
has
\[
 (\lambda_1,\ldots,\lambda_{2g-3})=s\eta,
 \qquad \lambda_{2g-2}>\frac14.
\]
\end{proposition}

\begin{proof}
Choose an invertible $m\times m$ spectral minor at $w^0$ and a degree domain
$D$ in those $m$ effective cuff variables, exactly as in
Proposition~\ref{prop:stable-transfer}.  The remaining rescaled cuff
weights and normalized twist angles give $4g-3$ complementary
Fenchel--Nielsen parameters.  Let them range in a sufficiently small compact
cell $Q$ in one augmented local chart.

The unselected rescaled cuff weights can change the limiting graph map.
Denote that map in the selected variables by $\Lambda_{\Gamma,q}$; twist
coordinates do not enter it.  At the base value $q^0$, the chosen degree
domain has degree $\pm1$ over a target ball.  After shrinking $Q$ and the
ball, the straight-line homotopy from $\Lambda_{\Gamma,q^0}$ to
$\Lambda_{\Gamma,q}$ avoids the target on $\partial D$, uniformly in
$q\in Q$.

The complementary pieces and all rescaled data range over a compact family
on $\overline D\times Q$.  The parameter-uniform form of
Lemma~\ref{lem:stable-chart-uniform} consequently applies there.  If
$\Phi_{t,q}$ is the scaled surface spectral map in the selected weights,
then
\[
 \sup_{(w,q)\in\overline D\times Q}
 \|\Phi_{t,q}(w)-\Lambda_{\Gamma,q}(w)\|\longrightarrow0
 \qquad(t\downarrow0).
\]
The displayed uniform convergence makes the surface-to-graph homotopy
boundary-avoiding simultaneously for every $q\in Q$.  Thus
\[
 \deg(\Phi_{t,q},D,\eta)
 =\deg(\Lambda_{\Gamma,q},D,\eta)
 =\deg(\Lambda_{\Gamma,q^0},D,\eta)=\pm1
\]
for every $q\in Q$, every target in the smaller ball, and every sufficiently
small $t$.  Taking $t=2\pi^2s$ proves the claim, and Otal--Rosas separates
the next eigenvalue.
\end{proof}

\begin{theorem}[Continuum maximal-prefix nonrigidity]
\label{thm:continuum-prefix}
The sets $V_g$ and constants $s_g$ in
Theorem~\ref{thm:ramanujan-wedge} may be chosen so that every
$s\eta$, $\eta\in V_g$ and $0<s<s_g$, is shared by continuum many pairwise
nonisometric closed genus-$g$ hyperbolic surfaces.  All these surfaces have
the same entire sub-$1/4$ spectrum, with multiplicity.

For each fixed $s$, almost every target in $sV_g$ admits a regular realization
in a marked Teichm\"uller chart.  Near such a realization its marked isolevel
is a real-analytic manifold of dimension $4g-3$; after passage to moduli the
corresponding local object is an orbifold quotient.
\end{theorem}

\begin{proof}
Apply Proposition~\ref{prop:parametric-transfer} to the regular Ramanujan
weighting in the proof of Theorem~\ref{thm:ramanujan-wedge}.  For every
$q\in Q$ the degree argument supplies at least one selected weight vector.
The resulting set of marked surfaces projects onto the continuum set $Q$.
Each unmarked isometry class has at most countably many markings, since the
mapping class group is countable.  Hence the family contains continuum many
unmarked isometry classes.  The $1/4$ statement follows from
Otal--Rosas and the choice $s\eta_m<1/4$.

On the locus where the prescribed eigenvalues are distinct and separated
from the next eigenvalue, the spectral map on marked Teichm\"uller charts is
real analytic.  Every preimage of a target in $sV_g$ lies in this locus:
the target entries are distinct and below $1/4$, while Otal--Rosas puts
$\lambda_{m+1}>1/4$.  A countable atlas covers this locus.  The union of the
critical-value sets of its chart maps is null by Sard's theorem, while the
parametric degree argument shows that every point of $sV_g$ has a preimage.
Thus almost every target has a regular marked realization, and the analytic
regular-level theorem gives dimension
$(6g-6)-(2g-3)=4g-3$.  Quotienting by the mapping class action gives the
stated local orbifold interpretation.  This is finite-prefix nonrigidity,
not full isospectrality.  It occurs in the degenerating thin regime, in
contrast with the finite-spectral-determination theorem of
Buser--Courtois on each fixed thick part \cite{BuserCourtois1990}.
\end{proof}

\begin{theorem}[Superexponentially many thin chambers]
\label{thm:many-chambers}
As $g\to\infty$, there are
\[
 \text{at least }\exp\{g\log g-O(g)\}
\]
pairwise nonisomorphic all-pants dual-graph chambers which share one exact
maximal sub-$1/4$ prefix.  More precisely, for every sufficiently large
$g$ there are a simple lacunary list $\rho_g$ and $s_g>0$ such that every
$s\rho_g$, $0<s<s_g$, is realized by continuum many surfaces in each of
that many chambers.
\end{theorem}

\begin{proof}
Put $N=2g-2$ and choose one simple list $\rho_g$ in the universal cone
$\cL_N(\delta_N)$.  Every connected simple cubic graph on $N$ vertices has
a strict regular realization of $\rho_g$ by
Theorem~\ref{thm:intro-universal}.  Such a graph has $3g-3$ edges and first
Betti number $g$, so it is an all-pants dual graph.  Apply
Proposition~\ref{prop:parametric-transfer} separately to every graph chart.
There are only finitely many types for fixed $g$, so the minimum of their
positive exact-lift thresholds is positive.  This repeated parametric-degree
argument, rather than an appeal to a single Ramanujan chart, gives continuum
many realizations inside every chamber.

The configuration-model enumeration of labeled simple cubic graphs is
\cite{Bollobas1980}
\[
 (e^{-2}+o(1))
 \frac{(3N)!}{(3N/2)!\,2^{3N/2}\,6^N}.
\]
Asymptotically almost every cubic graph is connected; for example this
follows from the stronger Hamiltonicity theorem
\cite{RobinsonWormald1992}.  Dividing by $N!$ gives a lower bound on the
number of unlabelled types whose logarithm is
\[
 \frac N2\log N+O(N)=g\log g+O(g).
\]
Finally, after decreasing the common scale threshold, the canonical short
multicurve recovers the dual graph.  Nonisomorphic dual graphs therefore
belong to distinct isometry classes of thin chambers.
\end{proof}

\section{Essential collision strata}\label{sec:essential}

At a repeated eigenvalue, ordered eigenvalues are not differentiable and the
preceding Jacobian degree cannot be used.  The appropriate object is the
low-energy quadratic form.  Let $V_0=\one^\perp\subset\R^N$.  An
\emph{all-pants dual graph} is a connected cubic pseudograph representing a
decomposition into genus-zero pairs of pants; its loop cuffs have zero graph
energy.

\begin{lemma}[Uniform effective-form convergence]
\label{lem:effective-forms}
Fix a strictly weighted genus-$g$ all-pants dual graph $\Gamma$ and a compact
set $C$ of positive effective weights, with the positive parallel-edge
splittings, loop coefficients, and compact auxiliary data fixed as in
Lemma~\ref{lem:stable-chart-uniform}.  Equip $V_0$ with
$\langle F,G\rangle_0=2\pi F^{\mathsf T}G$.  Let
$\mathcal E_{t,w}$ be the span of the first $m$ positive surface eigenspaces.
There is a resulting continuous isometry
$U_{t,w}:V_0\to\mathcal E_{t,w}$, obtained from collar-harmonic test
functions, low-energy spectral projection, and polar decomposition.  Define
the pulled-back operator $Q_{t,w}$ by
\[
 \langle Q_{t,w}F,G\rangle_0
 =\int_{X_\Gamma(t,w)}
   \langle\nabla U_{t,w}F,\nabla U_{t,w}G\rangle.
\]
Then
\begin{equation}\label{eq:effective-form-limit}
 \sup_{w\in C}
 \left\|\frac{2\pi^2}{t}Q_{t,w}
       -L_\Gamma(w)|_{V_0}\right\|\longrightarrow0
 \qquad(t\downarrow0).
\end{equation}
For isolated clusters, the corresponding Kato-transported block forms also
converge uniformly.
\end{lemma}

\begin{proof}
Let $J_{t,w}F$ equal the constant $F_v$ on the thick core of the component
indexed by $v$, interpolate harmonically across each collar, and subtract
its mean.  This is the standard test embedding in the graph-degeneration
argument of Colbois--Colin de Verdi\`ere
\cite[Lemma~V.1]{ColboisCdV1988}.  Uniformly for $w\in C$, collar capacity
and the area of a pair of pants give
\begin{align*}
 \langle J_{t,w}F,J_{t,w}G\rangle_{L^2}
   &=2\pi F^{\mathsf T}G+o(1)\|F\|\|G\|,\\
 \int\langle\nabla J_{t,w}F,\nabla J_{t,w}G\rangle
   &=\frac{t}{\pi}F^{\mathsf T}L_\Gamma(w)G
     +o(t)\|F\|\|G\|.
\end{align*}
Indeed, a collar of core length $tw_e$ has capacity
$tw_e/\pi+o(t)$, and all thick data range in a compact set.

Let
$\widehat J_{t,w}=J_{t,w}(J_{t,w}^{*}J_{t,w})^{-1/2}$ be the polar
normalization of the test embedding, and let $R_{t,w}$ be the operator on
$V_0$ obtained by pulling the Dirichlet form back through $\widehat J_{t,w}$.
The two displayed estimates give
\begin{equation}\label{eq:test-form-limit}
 \sup_{w\in C}\left\|
 R_{t,w}-\frac{t}{2\pi^2}L_\Gamma(w)|_{V_0}
 \right\|=o(t).
\end{equation}
The test space has dimension $m$ and Rayleigh quotient $O(t)$.  Hence the
min--max principle gives $m$ eigenvalues bounded by $O(t)$, while
Otal--Rosas gives $\lambda_{m+1}>1/4$.  If $P_{t,w}$ is the orthogonal
projection onto $\mathcal E_{t,w}$, the resulting uniform gap estimate makes
$B_{t,w}=P_{t,w}\widehat J_{t,w}$ invertible for small $t$; explicitly,
\[
 \|(I-P_{t,w})\widehat J_{t,w}F\|_{L^2}^2
 \le 4\int|\nabla\widehat J_{t,w}F|^2
 =O(t)\|F\|_0^2.
\]
Take the polar part
\[
 U_{t,w}=B_{t,w}(B_{t,w}^{*}B_{t,w})^{-1/2},
\]
where the adjoint uses $\langle\cdot,\cdot\rangle_0$ on $V_0$.
The stronger small-eigenvalue form comparison of
\cite[Lemma~V.1]{ColboisCdV1988}, applied uniformly on $C$, gives
\begin{equation}\label{eq:projected-test-form-error}
 \sup_{w\in C}\|Q_{t,w}-R_{t,w}\|=o(t).
\end{equation}
This is the step that controls the energy of the discarded high-mode
component; the preceding gap estimate alone only proves invertibility.
Combining \eqref{eq:test-form-limit} and
\eqref{eq:projected-test-form-error} proves
\eqref{eq:effective-form-limit}.  The contour formula for Riesz projections
then gives uniform convergence of every isolated Kato block
\cite{Kato1995}.
\end{proof}

Let $A=L_\Gamma(w^0)|_{V_0}$ have distinct eigenvalues
$\alpha_1<\cdots<\alpha_r$ and eigenspaces $E_i$.  Kato transport identifies
the nearby cluster spaces with $E_i$ and gives a block map
\[
 \mathcal Q_\Gamma(w)\in\bigoplus_{i=1}^r\Sym(E_i).
\]
Following the essential-family terminology, we say that the family has the
\emph{weak Arnold hypothesis} (WAH) at $w^0$ if the target
$\mathcal Q_\Gamma(w^0)$ lies in the interior of the image of every
sufficiently small continuous perturbation of this block map on some compact
neighborhood; compare the stable multiplicity framework of
\cite{CdVArnold1988}.  Surjectivity of
\begin{equation}\label{eq:SAH-derivative}
 D\mathcal Q_\Gamma(w^0):\R^{E_{\rm eff}(\Gamma)}
 \longrightarrow\bigoplus_i\Sym(E_i)
\end{equation}
is the strong Arnold hypothesis (SAH) and implies WAH.
Indeed, choose a linear complement $S$ to the kernel of the derivative in
\eqref{eq:SAH-derivative}.  On the affine slice $w^0+S$ the derivative is an
isomorphism.  A sufficiently small closed ball in this slice lies in the
positive cone and has nonzero local Brouwer degree; every sufficiently small
$C^0$ perturbation has the same boundary degree and therefore still hits the
target.

\begin{theorem}[Essential maximal-cusp transfer]
\label{thm:essential-transfer}
If a trace-one graph direction $p$ has a strict WAH realization on a
genus-$g$ all-pants dual graph, then there is $s_p>0$ such that every $sp$,
$0<s<s_p$, is exactly the first $m$ positive-eigenvalue list, with all
multiplicities preserved.
\end{theorem}

\begin{proof}
Choose the compact neighborhood in the WAH definition inside the positive
weight cone and shrink it if necessary.  By
Lemma~\ref{lem:effective-forms}, the rescaled surface block map is uniformly
$C^0$-close there to the graph block map for every sufficiently small $t$.
WAH therefore supplies a positive nearby weighting at which the surface
block form hits the target blocks exactly.  The eigenvalues of the rescaled
low operator $(2\pi^2/t)Q_{t,w}$, including all multiplicities, are therefore
$p$, so the actual surface eigenvalues are $(t/(2\pi^2))p$.  Putting
$t=2\pi^2s$ gives every sufficiently small multiple $sp$; the pointwise
Otal--Rosas theorem places the next eigenvalue above $1/4$.
\end{proof}

The SSPWL is sufficient for the hypothesis above.  Indeed,
\cite[Proposition~2.9]{CFGL2026} identifies SSPWL with transversality of the
graph-pattern space to the orthogonal isospectral orbit.  The normal quotient
to that orbit is $\bigoplus_i\Sym(E_i)$, so the projected derivative is
exactly \eqref{eq:SAH-derivative}.  Thus SSPWL implies SAH and WAH.  This
does not say that the lifted surface itself has SAH.

If the positive multiplicities are $r_1,\ldots,r_k$, a useful necessary
condition follows by counting dimensions:
\begin{equation}\label{eq:multiplicity-budget}
 \sum_i\frac{r_i(r_i+1)}2
 =m+\sum_i\binom{r_i}{2}\le3g-3,
 \qquad\text{hence}\qquad
 \boxed{\sum_i\binom{r_i}{2}\le g.}
\end{equation}
Loops and parallel cuffs can only reduce the number of independent effective
forms, so \eqref{eq:multiplicity-budget} is not sufficient.

\section{The sharp cusp variance and forbidden cone}
\label{sec:cusp-gap}

Let $b_m=(1/m,\ldots,1/m)$ and $V(p)=\|p-b_m\|_2^2$.  For a trace-one
effective graph write $q$ for the number of positive distinct nonloop pairs,
$d_v=\sum_{e\ni v}w_e$, and $W=\sum_e w_e=1/2$.  Then
\begin{equation}\label{eq:trace-square}
 \tr L^2=\sum_vd_v^2+2\sum_ew_e^2
 \ge\frac1N+\frac1{2q}\ge\frac4{3N},
\end{equation}
because $q\le3N/2$.  This includes disconnected positive supports and hence
the multiscale boundary.

More precisely, \eqref{eq:trace-square} has the following exact quadratic
SOS form:
\begin{equation}\label{eq:sharp-SOS}
 \boxed{
 \tr L^2-\frac4{3N}
 =\sum_v\left(d_v-\frac1N\right)^2
  +2\sum_{e=1}^q\left(w_e-\frac1{3N}\right)^2
  +\frac{3N-2q}{9N^2}.}
\end{equation}
The last slack is essential unless $q=3N/2$.  Equality occurs precisely for
a loopless simple cubic effective support with all weights $1/(3N)$.
Weighted-Laplacian minimum-variance inequalities and the equal-weight
regular-graph minimizer already appear in \cite[Section~4]{FGL2025}.  The
point here is the optimization over every stable genus-$g$ type, including
zero-weight faces, and its sharp transfer to the surface limit-inf.

\begin{theorem}[Sharp maximal-cusp variance]\label{thm:sharp-cusp}
For $g\ge3$, with $N=2g-2$ and $m=N-1$,
\begin{equation}\label{eq:sharp-cusp-limit}
 \boxed{
 \lim_{\varepsilon\downarrow0}
 \inf_{\lambda_m(X)\le\varepsilon}V(p(X))
 =c_g:=\frac{N-4}{3N(N-1)}
 =\frac{g-3}{3(g-1)(2g-3)}.}
\end{equation}
For every $g\ge4$ and every $r<\sqrt{c_g}$, there is $\varepsilon_r>0$
such that
\[
 \lambda_m(X)\le\varepsilon_r
 \quad\Longrightarrow\quad
 \|p(X)-b_m\|_2\ge r.
\]
The radius $\sqrt{c_g}$ is sharp at the tangent-link level.
\end{theorem}

\begin{proof}
For $p\in\mathcal T_g$, trace normalization gives
$\sum_jp_j^2=\tr L^2$.  Hence \eqref{eq:trace-square} yields
\[
 V(p)=\sum_jp_j^2-\frac1{N-1}
 \ge\frac4{3N}-\frac1{N-1}=c_g.
\]
The bound is attained by the equal weighting of $K_4$ when $N=4$ and, for
every even $N\ge6$, by an equal-weight connected simple cubic graph, for
example the prism $C_{N/2}\mathbin\square K_2$.  Such a graph is an all-pants
dual graph of genus $N/2+1=g$, so Theorem~\ref{thm:tangent-link} makes the
bound sharp on the link.

To pass from the link to \eqref{eq:sharp-cusp-limit}, take an arbitrary
sequence with $\lambda_m\to0$ and then a convergent subsequence of its
normalized directions; every limit lies in $\mathcal T_g$.  Conversely, a
surface sequence approaches every link point by definition.  Compactness of
the link gives equality of the limit-inf and $\min_{\mathcal T_g}V$.  The
finite-scale exclusion for every smaller radius follows immediately.
\end{proof}

For $g=2$ the link consists of one normalized coordinate and $V=0$; the
formula for $c_g$ is not used.  For $N>4$, an equal-weight simple cubic graph
with simple Laplacian spectrum supplies a genuine skeleton point.  Indeed,
its direction minimizes $V$ globally on $\mathcal T_g$.  A regular preimage
would make that direction interior, but moving slightly toward $b_m$ lowers
$V$ below $c_g$.  Thus no stable graph realization of that direction is
normalized regular.  The skeleton is therefore not merely a defect of the
Sard argument.

\begin{corollary}[A cusp condition-number bracket]
\label{cor:cusp-condition-bracket}
For $g\ge3$, put $N=2g-2$, $m=N-1$, and
\[
 \kappa_g^*=\inf_{p\in\mathcal U_g}\frac{p_m}{p_1},
 \qquad
 q_g=\frac{N-4}{3N}.
\]
Then
\begin{equation}\label{eq:cusp-condition-bracket}
 \left(\sqrt{1+q_g}+\sqrt{q_g}\right)^2
 \le\kappa_g^*<18+12\sqrt2.
\end{equation}
Equivalently, the lower bound is
\[
 \frac{5N-8+4\sqrt{(N-1)(N-4)}}{3N}.
\]
It tends to $3$ as $g\to\infty$.  The bounds do not determine the optimum.
\end{corollary}

\begin{proof}
Let $p\in\mathcal T_g$ be positive and put $R=p_m/p_1$.  For numbers
$x\in[a,b]$ with mean $\mu$, averaging $(x-a)(b-x)\ge0$ and optimizing the
result over $\mu/a\in[1,b/a]$ gives
\[
 \frac{\operatorname{mean}(x^2)}{\mu^2}
 \le\frac{(R+1)^2}{4R}.
\]
Apply this to the $m$ coordinates of $p$.  Since $\sum_jp_j=1$,
\[
 1+mV(p)=m\sum_jp_j^2\le\frac{(R+1)^2}{4R},
\]
and hence
\[
 mV(p)\le\frac{(R-1)^2}{4R}.
\]
Theorem~\ref{thm:sharp-cusp} gives
$mV(p)\ge mc_g=q_g$.  Inverting the last inequality on the branch $R\ge1$
gives the left side of \eqref{eq:cusp-condition-bracket}.  The regular
Ramanujan direction constructed in Theorem~\ref{thm:ramanujan-wedge} gives
the strict upper bound.  The alternative formula and the limit are direct
algebra.  This corollary is a range reformulation of the sharp variance
bound, not a separate optimality theorem.
\end{proof}

\section{The classical genus-three prescription as endpoint calibration}
\label{sec:genus-three}

The vanishing of $c_3$ reflects exact flexibility rather than only failure of
the obstruction.

\begin{theorem}[Genus-three prescription]\label{thm:genus-three}
For every $0<a\le b\le c$ there is $s_0>0$ such that, for every
$0<s<s_0$, a closed genus-$3$ hyperbolic surface satisfies
\[
 (\lambda_1,\lambda_2,\lambda_3)=s(a,b,c),
 \qquad \lambda_4>\frac14.
\]
Repeated entries are preserved exactly.
\end{theorem}

\begin{proof}
On $K_4$ put
\[
 z=\frac a4,\qquad y=\frac{b-z}{3},\qquad
 x=\frac{c-y-z}{2},
\]
and assign
\[
 w_{12}=x,\qquad w_{13}=w_{23}=y,\qquad
 w_{14}=w_{24}=w_{34}=z.
\]
The inequalities $0<a\le b\le c$ give $x,y,z>0$.  The mutually orthogonal
vectors of types $(1,1,1,-3)$, $(1,1,-2,0)$, and $(1,-1,0,0)$ show that the
positive spectrum is
\[
 (4z,3y+z,2x+y+z)=(a,b,c).
\]

Moreover, the linear map
\begin{equation}\label{eq:K4-form-isomorphism}
 \R^{E(K_4)}\longrightarrow\Sym(\one^\perp),\qquad
 w\longmapsto L_{K_4}(w)|_{\one^\perp},
\end{equation}
is an isomorphism between six-dimensional spaces.  On functions modulo
constants use coordinates $r_i=f_i-f_4$, $1\le i\le3$.  The three edge forms
for $i4$ are $r_i^2$, while those for $ij$ are $(r_i-r_j)^2$; together they
span every quadratic form in $(r_1,r_2,r_3)$.  Thus this strict family has
SAH, including at collisions.

Put $S=a+b+c$ and $p=(a,b,c)/S$, and divide the displayed edge weights by
$S$.  This gives the trace-one direction required by
Theorem~\ref{thm:essential-transfer}.  That theorem produces the exact
surface form for every sufficiently small multiple of $p$; reparameterizing
the scale by $S$ gives every sufficiently small $s(a,b,c)$.  Form equality,
rather than differentiability of ordered eigenvalues, preserves the
multiplicities.  Finally use Otal--Rosas for $\lambda_4$.
\end{proof}

The classical antecedent is property $A_4$ of Colin de Verdi\`ere
\cite[Theorem~1.3 and \S2a]{CdV1987}: in genus three it prescribes the zero
mode and three positive eigenvalues, with multiplicities, using metrics that
may be chosen of constant curvature.  Because its curvature/area
normalization is not fixed, it is not by itself the curvature-$-1$,
all-sufficiently-small-scale statement above.  Our proof is a direct $K_4$
normalization of the classical essential-family mechanism.  We claim no
novelty for that mechanism; the contribution used elsewhere in the paper is
the fixed-cusp formulation together with the surrounding tangent-link and
sharp genus-$3/4$ transition.

\section{Semialgebraic complexity and SOS certificates}
\label{sec:complexity}

For rational ordered $p=(p_1,\ldots,p_m)$ with $p_j\ge0$ and
$\sum_jp_j=1$, define
\[
 \textsf{MCTL}=\{(g,p):p\in\mathcal T_g\}.
\]

\begin{theorem}[Complexity upper bounds]\label{thm:MCTL-complexity}
One has
\[
 \textsf{MCTL}\in\exists\mathbb R\subseteq\mathrm{PSPACE}.
\]
For every fixed $g$, rational membership in $\mathcal T_g$ is decidable in
deterministic time polynomial in the total bit length of $p$; the polynomial
and its constants may depend on $g$.
\end{theorem}

\begin{proof}
It is not enough, when $g$ is part of the input, to enumerate all stable
graphs.  Instead give each of the $N$ labeled vertices three half-edges and
use Boolean real variables $x_{\{a,b\}}$, indexed by unordered pairs of
distinct half-edges, to encode a perfect matching:
\[
 x_{\{a,b\}}^2-x_{\{a,b\}}=0,\qquad
 \sum_{b\ne a}x_{\{a,b\}}=1.
\]
A standard single-commodity flow of polynomial size enforces connectedness
of the matched pseudograph.  Give a matched nonloop pair a variable
$0\le w_{\{a,b\}}\le x_{\{a,b\}}/2$, set loop energy to zero, and impose
$2\sum w_{\{a,b\}}=1$ over nonloop pairs.  Parallel terms simply add in the resulting Laplacian
$L$.  A connected cubic pseudograph on $N$ vertices has first Betti number
$3N/2-N+1=g$.  A stable vertex of type $(h,\deg)=(1,1)$ is represented by a
genus-zero cubic vertex whose two additional half-edges are paired into a
self-loop; this changes neither the effective graph energy nor its spectral
image.  Thus the cubic encoding covers every type in $\mathfrak G_g$, and
zero conductances recover all of their weight faces.

Introduce $U\in\R^{N\times m}$ and impose
\begin{equation}\label{eq:ETR-eigenbasis}
 U^{\mathsf T}U=I_m,\qquad U^{\mathsf T}\one=0,\qquad
 L=U\diag(p_1,\ldots,p_m)U^{\mathsf T}.
\end{equation}
These are polynomially many quadratic equations and inequalities, giving a
polynomial-size existential-real formula.  Canny's theorem yields the PSPACE
upper bound \cite{Canny1988}.  When $g$ is fixed, the real dimension, degrees,
and number of constraints are fixed, so Renegar's bounds give polynomial bit
complexity \cite{Renegar1992I}.  No hardness or uniform fixed-parameter claim
is made.
\end{proof}

Regular membership also has a polynomial existential witness.  If
$J_\Gamma(w)$ is the $m$-row unnormalized coefficient Jacobian, use a right
inverse variable
\begin{equation}\label{eq:right-inverse-witness}
 J_\Gamma(w)R=I_m
\end{equation}
instead of expanding exponentially many minors.

There is a complementary graph-by-graph certificate for nonmembership.  For
a fixed effective type with $q$ edges put
\[
 K_\Gamma=\{w\ge0:2\sum_ew_e=1\}
\]
and, for trace-one $p$, define
\begin{equation}\label{eq:moment-objective}
 F_{\Gamma,p}(w)=
 \sum_{k=2}^{N}\left(
    \tr L_\Gamma(w)^k-\sum_{j=1}^{m}p_j^k
 \right)^2.
\end{equation}
The trace fixes the first power sum, and Newton identities show
\begin{equation}\label{eq:moment-membership}
 p\in\psi_\Gamma(K_\Gamma)
 \quad\Longleftrightarrow\quad
 \min_{K_\Gamma}F_{\Gamma,p}=0.
\end{equation}

If $p$ is not in this graph image, compactness makes the minimum
$\delta_{\Gamma,p}>0$.  Add the redundant ball inequality
$b(w)=1-\|w\|^2\ge0$ and let $h(w)=1-2\sum_ew_e$.  The resulting quadratic
module is Archimedean.  For $0<\epsilon<\delta_{\Gamma,p}$, Putinar's
Positivstellensatz gives, at some finite degree,
\begin{equation}\label{eq:putinar-certificate}
 F_{\Gamma,p}-\epsilon
 =\sigma_0+\sigma_bb+\sum_e\sigma_ew_e+\tau h,
\end{equation}
where the $\sigma$'s are sums of squares and $\tau$ is arbitrary
\cite{Putinar1993}.  Equivalently, some finite Lasserre relaxation separates
the optimum from zero \cite{Lasserre2001}.  Since the stable types are finite
for fixed $g$, a nonmember of $\mathcal T_g$ has one such certificate for
every type.

This conclusion has three important limits.  There is no uniform bound on
the SOS degree or bit size; the coefficients in
\eqref{eq:putinar-certificate} may be real; and membership or boundary
nonnegativity is not certified by the same strict-positivity argument.  In
particular, this does not place the complement in $\mathrm{NP}$ or make the
Lasserre hierarchy a two-sided decision procedure.

\section{Residual directions and the global problem}\label{sec:residual}

The preceding theorems leave a lower-dimensional but genuine residual set.
Two cubic graphs make the boundary of the argument explicit.

\subsection{The genus-four collision barrier}

The unit-weight Laplacian of $K_{3,3}$ has spectrum
\[
 (0,3,3,3,3,6),\qquad
 p_{3,3}=\left(\frac16,\frac16,\frac16,\frac16,\frac13\right),
 \qquad V(p_{3,3})=c_4=\frac1{45}.
\]
Its positive multiplicities are $4+1$, so the block-form target dimension
is $10+1=11$, whereas a genus-$4$ pants family has at most $3g-3=9$
independent cuff parameters.  Hence no leading graph realization of this
pattern can satisfy WAH or SAH.  This is only a first-order barrier: it proves
neither realization nor nonrealization of the exact surface ray.

\subsection{The genus-seven Frucht fold}

Let $F$ be the Frucht graph, with $12$ vertices and $18$ edges.  Its
adjacency characteristic polynomial factors as
\begin{equation}\label{eq:frucht-factorization}
 (x-3)(x-2)x(x+1)(x+2)
 (x^3+x^2-2x-1)(x^4+x^3-6x^2-5x+4).
\end{equation}
The two nonlinear factors have discriminants $49$ and $22545$, resultant
$64$, and no roots among $3,2,0,-1,-2$.  Thus the equal-weight Laplacian has
simple positive spectrum.  Its normalized direction $p_F$ satisfies
$V(p_F)=c_7=2/99$ and is a global minimizer of $V$ on $\mathcal T_7$,
hence a critical value of the normalized graph map.

\begin{proposition}[Certified definite Frucht fold]\label{prop:frucht-fold}
At unit weights, the unnormalized coefficient Jacobian of $F$ has rank $10$
and contains a minor of determinant $73728$.  On the fixed-total-weight
simplex, the restriction of the coefficient Jacobian has rank $9$ and
contains a minor of determinant $36864$.  Consequently the normalized
spectral derivative has rank $9$, and the normalized spectral map has a
corank-one definite fold at $p_F$.
\end{proposition}

\begin{proof}
The accompanying exact-arithmetic certificate computes the characteristic
coefficients through the multiaffine weighted-forest formula and verifies the
two integer minors.  The variance boundary supplies the matching upper ranks
$10$ and $9$, so the computed ranks are exact.  On the fixed-weight-sum
slice let $F_0$ be the normalized spectral map and $h=V-c_7$.  Then
$dh(p_F)=2(p_F-b_m)$ spans the one-dimensional cokernel.  The cubic equality
case of \eqref{eq:sharp-SOS} gives, with
$W=\sum_e w_e$, $\bar d=2W/12$, and $\bar w=W/18$,
\[
 V(F_0(w))-c_7=
 \frac{\sum_v(d_v-\bar d)^2+2\sum_e(w_e-\bar w)^2}{4W^2},
\]
whose Hessian is positive definite on the eight-dimensional kernel of
$DF_0$.  The Morse lemma with parameters gives local coordinates
\[
 (x,y)\longmapsto(x,\|y\|^2),
 \qquad x\in\R^9,\quad y\in\R^8.
\]
The underlying cubic graph appears in the catalogue
\cite{BussemakerCobeljicCvetkovicSeidel1976}; the displayed factorization and
both minors are independently verified by the accompanying exact-arithmetic
certificate.  The definiteness of the fold is supplied analytically by the
displayed SOS identity, not by the rank certificate alone.
\end{proof}

This proposition is a graph theorem, not an exact surface-ray theorem.
Burger's $C^0$ comparison does not determine which side of a fold the
surface perturbation occupies.  If a fixed plumbing family admits a refined
uniform expansion
\[
 F_t=F_0+a(t)R+o_{C^2}(a(t)),\qquad a(t)\downarrow0,
\]
Lyapunov--Schmidt reduction leaves a scalar normal equation
\[
 \|y\|^2+a(t)\beta_F+o(\|y\|^2+a(t))=0.
\]
The sign of the still-uncomputed coefficient $\beta_F$ decides local
existence in that family.  It does not by itself give a global exclusion
over all twists and graph-limit representations.

\subsection{Open ends}

The exact link, generic simple directions, and WAH collision directions do
not solve the finite-scale problem of \cite[Question~1]{EJT2026}.  Among the
remaining concrete problems are:
\begin{enumerate}[label=\textup{(\arabic*)},leftmargin=2.2em]
\item decide the $K_{3,3}$ and Frucht surface rays, the latter by computing
the signed normal plumbing correction;
\item build a compactified spectral cobordism controlling nonseparating
Deligne--Mumford escape, eigenvalue collisions, and crossings at $1/4$;
\item determine the optimal universal aperture $\delta_n^*$ and the true
asymptotic order of the path volume $P(P_n)$.
\item determine the exact regular cusp condition number $\kappa_g^*$ between
the bounds in Corollary~\ref{cor:cusp-condition-bracket}, and describe its
minimizing graph types.
\end{enumerate}
Thus the maximal-cusp tangent link is completely characterized, and its
exact surface germ is settled off a compact semialgebraic null skeleton and
on all essential collision strata.  Inside the regular germ there is,
for every genus, a genus-dependent full-dimensional wedge whose internal
condition number is bounded by the genus-independent constant
$18+12\sqrt2$ and whose every target has continuum many realizations;
common lacunary targets also persist across
superexponentially many all-pants thin regions.  The residual skeleton,
optimal condition number, and global inverse problem remain open.

\bibliographystyle{amsplain}
\bibliography{references}

\end{document}
